\documentclass[10pt]{article}
\usepackage{tmlr}
\usepackage{amsmath,amssymb,booktabs,longtable,array,calc}
\usepackage{graphicx,xcolor,tikz}
\usetikzlibrary{arrows.meta,positioning}
\usepackage{hyperref,url}
\newcommand{\tightlist}{\setlength{\itemsep}{0pt}\setlength{\parskip}{0pt}}
\hypersetup{hidelinks}

\title{What Decision Records Establish About\\Model Choices Under Uncertainty}
\author{Anonymous authors}
\begin{document}
\raggedbottom
\maketitle
\begingroup\renewcommand{\thefootnote}{}\footnotetext{AI assistance: An LLM-based assistant contributed to experimental design, implementation, analysis, literature review, and manuscript drafting and revision. AI tools are not authors.}\endgroup
\begin{abstract}
Decision records can be evaluated by whether they match an action,
identify an executable selection rule, or help predict a later choice.
We distinguish these questions through a case study of model decisions
under uncertainty. In a concrete counterexample, five ways of completing
an incompletely specified selector reproduce the same original action:
one plan is best on all five criteria. Applied to another fixed
representation, those selectors yield different preservation/loss labels
under the study's conditional feasibility evaluation. Original-action
agreement therefore leaves the resulting component diagnosis ambiguous.
In a separate follow-up, required records match all sixteen model proposals
executed without external selection. The two supported rules agree in all
sixteen records, with fourteen agreements forced by a single row and two
observed in multirow tables. Finally, a prospective pilot freezes eight new
records before sixteen target choices across two execution segments.
Prediction with versus without the record is 6/8 versus 7/8 for Astra and
3/8 versus 3/8 for Sol: no positive net accuracy gain is observed across
four purposeful histories per profile. The study provides concrete controls
for evaluating decision records: test whether candidate rules can disagree,
separate action correspondence from task labels, and assess added prediction
on later choices. These measurements concern external records in two task
families; population effects require broader sampling.
\end{abstract}
\section{Question and contribution}
Our Research Under Uncertainty (RUU) program studies how models form assumptions, choose
experiments and revise predictions when an environment's mechanism is unknown.
Decision records offer a way to examine those activities. Before using them
as evidence, however, we need to know which relationships they measure.
This paper asks: \textbf{when does a record that matches a model's action
identify a selection rule, and when does it help predict a later choice?}

We study externally supplied tables and descriptions in two simulated task
families. An executable record combines an action-scoring representation
with a deterministic rule for selecting an action. For an incompletely
specified record, the analyst must supply the remaining choices; we call
this a \emph{selector completion}. Combining one record's representation
with another record's rule is a \emph{crossing}. For example, applying
column means and column maxima to the same table can select different
actions. Our initial component diagnostic used such crossings to examine
which part of a record supported a successful choice.

\paragraph{A concrete identification problem.}
In the pilot's original representation A2, one plan is best on all five
examined criteria. Total cost and trajectory peak therefore reproduce the
same model action. In a second representation A1, held fixed for comparison,
the criteria separate: total cost selects a plan with a preservation witness,
while trajectory peak selects a different plan with a conditional loss
certificate. Agreement on A2 leaves these completions indistinguishable,
although they give different crossed diagnoses
(Section~\ref{sec:selector-sensitivity}). This example supplies a practical
check: \emph{does the validation situation give candidate selectors an
opportunity to disagree?} When it does not, the diagnosis must remain
conditional on the independently specified selection contract.

\paragraph{Correspondence and prediction are separately measurable.}
A follow-up separates requiring a record from externally selecting its
action. Each profile supplies eight records matching its own executed
proposal. Fourteen of these sixteen numeric tables have one row, so both
supported rules necessarily agree; their agreement in the remaining two
multirow tables is observed. We pair these counts with an action-copy control
and a scan over the supported rules (Section~\ref{sec:followup}). This makes
both correspondence and its opportunity for rule discrimination visible.

A prospective pilot then tests added prediction on new choices. Records are
frozen before two target branches per source, and paired observers receive
the same information with or without the record. The observed net accuracy
differences are $-1/8$ for Astra and $0/8$ for Sol across four purposeful
histories per profile (Section~\ref{sec:prospective}). The pilot measures
transfer from a source that reports an action and record to a target that
reports an action alone; it evaluates the numeric table and accompanying
prose together.

\paragraph{Contribution to evaluation practice.}
The evaluation decision at issue is accepting an analyst-completed selector
because it reproduces the source action. The archived A2 case exposes the
problem: one plan is best on all five criteria, so every tested completion
passes that check, yet the same completions produce different choices and
conditional task diagnoses on A1. Researchers using such records should
therefore retain the compatible selector set and test a situation where its
members can disagree before attributing a diagnosis to a particular rule.

The follow-up and prospective panels make two further reporting decisions
concrete: separate a record's match to an unforced proposal from the task
label, and measure added prediction on new choices after freezing the record.
An action change can preserve its P or L label
(Table~\ref{tab:followup-rule-examples}); an aggregate prediction difference
can combine opposite observed contrasts across tasks
(Section~\ref{sec:prospective}). These cases give researchers specific checks
to add to a record evaluation. The common-menu recovery identity and
single-row rule agreement explain what the contract already guarantees
(Appendix~\ref{app:structural-control}); the archived cases and model decisions
show where the distinctions matter in practice. Table~\ref{tab:study-map}
maps each panel to its question, and Section~\ref{sec:related} relates these
external-record measurements to explanation evaluation.

\begin{table}[htbp]
\centering\small
\begin{tabular}{@{}>{\raggedright\arraybackslash}p{0.14\linewidth}>{\raggedright\arraybackslash}p{0.23\linewidth}>{\raggedright\arraybackslash}p{0.25\linewidth}>{\raggedright\arraybackslash}p{0.29\linewidth}@{}}\toprule
Study & Question & Intervention or check & Supported claim\\\midrule
Pilot (8 trials) & Does native agreement identify a selector? & Five analyst completions of archived records; component crossings & Concrete non-identifiability under these completions\\
Transfer (48 trials) & Are free-form records crossing-ready? & Optional records across two task families & Applicability is limited; zero eligible pairs is not success\\
Common menu (32 trials) & What does the fixed numeric contract add? & Optional records versus required records with external selection & Structural recoverability; record requirement and execution remain combined\\
Follow-up (48 assignments) & Can record and proposal match without forced execution? & A optional/model, B required/model, C required/selector & Profile-specific correspondence, record structure and conditional labels; no conflicting C proposal observed\\
Prospective choice prediction (16 targets) & Does a frozen record add prediction value? & Paired observers with/without record; same public evidence and profile & No positive net accuracy gain observed; four histories/profile, two linked targets/history\\\bottomrule
\end{tabular}
\caption{Study map. Counts refer to distinct panels, not pooled replications.
The follow-up has 46 completed assignments and two retained technical losses.
The prospective pilot has 15 original targets plus one separately completed
target; its original interrupted request remains recorded.
The separate finite-class companion is not needed for the main argument.}
\label{tab:study-map}
\end{table}

\section{Task, record, execution and evaluation}
\label{sec:measurement}
\paragraph{Task.}
The actor receives an action/observation history, a remaining objective and a
finite menu, without the hidden mechanism. In the continuous task it chooses
a bounded two-dimensional control to preserve the possibility of reaching
and maintaining a target region under a path constraint. In the route task
it chooses a legal discrete action under remaining route constraints.
The common-menu and prospective panels offer 13 continuous controls or
three route actions, and allow actors up to three analyses of the observed
history. Both families supply the past actions and observed states, the
remaining objective and constraints; the hidden transition mechanism is
withheld. The prospective panel's concrete source--target changes appear in
Table~\ref{tab:prospective-branches}.

\paragraph{Record.}
The actor can submit a numeric table of candidate failure forecasts and a
declared rule. Columns identify menu actions; rows contain values in $[0,1]$.
Uniform mean minimizes the column mean; worst case minimizes the column
maximum; ties use the lowest index. The executable contract specifies action selection; forecast calibration
and the semantic adequacy of rows require separate evaluation.
Table~\ref{tab:real-record} shows one complete numeric table from the follow-up,
before any formal crossing notation is needed.

% Exact numeric values from finals/trial-003.json; no rounding.
\begin{table}[htbp]
\centering\small
\begin{tabular}{lrrr}\toprule
Archived candidate-loss row & Menu action 0 & Menu action 1 & Menu action 2\\\midrule
Only row & 0.35 & 0.30 & 0.25\\\bottomrule
\end{tabular}
\caption{A real B record: Astra, route history \texttt{probe-03}, repeat 1,
trial 3. The declared mean rule selects action 2, as does worst case.
The model proposed action 2 and B executed that proposal. Its bound menu label
is P. The numbers are uncalibrated forecasts; their ordering and the matching
proposal are reproducible, but their semantic adequacy or causal role is not
established. Exact final JSON and identity bindings are in the bundle.}
\label{tab:real-record}
\end{table}


\paragraph{Execution.}
The follow-up separates record requirement from action source
(Table~\ref{tab:arms-early}). In B the model proposal is executed even if a
valid record would select a different action. In C the external selector
determines execution, while the proposal remains observable. Thus B's
record--proposal agreement is observed; C's record--execution agreement is
guaranteed. Exact instructions and schema are in the bundle.

\begin{table}[htbp]
\centering\small
\begin{tabular}{lp{0.22\linewidth}p{0.23\linewidth}p{0.34\linewidth}}\toprule
Arm & Record requirement & Executed action & What agreement can establish\\\midrule
A & Optional & Model proposal & Correspondence, if a valid record exists\\
B & Required valid record & Model proposal & Observed record--proposal correspondence\\
C & Required valid record & Declared record selector & Execution match by construction; proposal match remains observable\\\bottomrule
\end{tabular}
\caption{Follow-up conditions. A valid record need not match the proposal to
be accepted in B. A missing record leaves correspondence undefined.}
\label{tab:arms-early}
\end{table}

\paragraph{Evaluation.}
The analyst attaches a task label to the executed menu action: preservation
(P), loss (L), or unresolved (U), using the bound feasibility evaluation.
That label is separate from whether a record is valid, matches the proposal,
or changes action under rule substitution. The model does not receive the
privileged evaluator during collection.

\subsection{Why matching needs a control}
For any proposed index $a$ in an $n$-action menu, construct one row
$A_{1j}=0$ when $j=a$, and $A_{1j}=1$ otherwise. Both supported rules select
$a$. This \emph{action-copy negative control} requires only the menu size
and proposed index, with no history, mechanism or task label. Conditional on
the proposal it supplies no independent information about how that proposal
was produced. It nevertheless satisfies validity and correspondence checks.
Consequently those checks alone cannot establish an informative explanation.
The construction tests the sufficiency of those checks. How the participant
produced its record and whether its forecasts are informative remain empirical
questions.

A positive arithmetic control has two rows,
$A=\left(\begin{smallmatrix}0&0.4\\0.7&0.4\end{smallmatrix}\right)$.
Mean scores $(0.35,0.4)$ select action 0; worst-case scores $(0.7,0.4)$
select action 1. The same reader therefore detects a rule contrast when the
values support one. Neither constructed control has a physical task label
or constitutes an additional model trial. The distributed check tests the
copy construction at every proposal index for menu sizes 3 and 13, under
each declared rule, and tests the positive control under both rules.

\subsection{Two external comparisons}
Let $A_i$ be record $i$'s table over a common finite menu and let $r_i$
be its declared deterministic rule. The action $a_i=r_i(A_i)$ is one
column of that menu. Here, $a_i$ denotes the record-selected action; the
model-proposed and executed actions are separate quantities. The
\emph{observed-rule crossing} has four entries
$c_{ij}=r_j(A_i)$ for $i,j\in\{1,2\}$. The separate
\emph{supported-rule-family table} evaluates each $A_i$ under both allowed
rules: uniform mean and worst case, with lowest-index tie breaking.
The family table may contain a rule that neither paired actor used. Its
entries are deterministic evaluations of archived records, with task labels
from the bound menu census; a live model counterfactual would require a new
model decision.

\subsection{Feasibility labels and access}
P requires a successful admissible continuation witness after the action.
L requires the specified loss certificate; absence of a witness alone is not
L. U preserves unresolved numerical cases. Continuous-control labels are
conditional on the disclosed floating-point and bounding premises in the
supplement. The discrete route family's finite computation has a different
verification contract. Both actors and analyst see the legal histories; only
the analyst receives the hidden mechanisms and privileged evaluator during
the original experiment. Reviewer disclosure now permits inspection of these
mechanisms without retroactively changing what participants received.

Record validity, correspondence and rule sensitivity can be checked without
the physical label map. P/L/U comparisons add its assumptions. The exact
arithmetic helper does not prove the uniform library bounds or attest the
historical runtime; those premises remain open. These are distinct evidence
levels when interpreting outcome differences.

\section{Evidence sets and their denominators}
The first three evidence sets below use crossing-pair denominators.
The separate follow-up in Section~\ref{sec:followup} instead compares A/B/C
conditions within eight matched triplets per profile; these condition pairs
are not crossing pairs. Its 48 assignments include two technical losses.
\begin{table}[htbp]
\centering\small
\begin{tabular}{p{0.28\linewidth}rrrr}
\toprule
Evidence set & Trials & Pairs & Eligible records & Eligible pairs\\
\midrule
Pilot (component study) & 8 & 4 & \multicolumn{2}{c}{2 selected component pairs}\\
Prospective transfer & 48 & 24 & 1 & 0\\
Optional record & 16 & 8 & 1 & 0\\
Mandatory record + execution & 16 & 8 & 16 & 8\\
\bottomrule
\end{tabular}
\caption{The three predecessor evidence sets. Pilot component eligibility is a different
construction from the frozen prospective adapter and common-menu contract.
Empty pair denominators leave the crossing effect unassessed.}
\end{table}

The pilot used four fixed histories with two repeats under the requested
gpt-6-astra/max profile. Its component study
focused on two boundary histories, pilot contexts B and D, after outcomes
were available. These history names are unrelated to the follow-up arms A/B/C.
The transfer panel used 48 live trials: two families, three instances per
family, two histories per instance, the gpt-6-astra/max and
gpt-6.1-sol/max profiles, and two repeats. Under its frozen adapter, a record had to
support the prescribed operation without post hoc completion. Only one trial
qualified, leaving zero of 24 pairs crossable. This measures the adapter's applicability to the available external
records. With no eligible pair, the crossing effect is unassessed; internal
representations are outside the adapter's observations. Native transfer labels were 42 P and 6 L. The adapter's mutually exclusive
outcomes were: 24 route trials with no analysis attempt, 18 whose last successful
analysis lacked a table, three whose rule did not reproduce the action, two
with a menu-identity mismatch, and one supported record. These reasons
distinguish absent evidence from an invalid submitted record.
The other 24, continuous-task trials made 72 analysis attempts: 67 succeeded
and five returned errors. The accounting distinguishes an unused analysis opportunity from a failed
analysis call. It measures visible tool use, leaving internal reasoning
outside the observation.

The 32-trial panel crossed two task families, two purposefully selected
histories per family, two model profiles (gpt-6-astra/max and gpt-6.1-sol/max),
two repeats and two conditions. Each condition has 16 trials and eight pairs.
A crossing pair comprises the two repeats at the same history and profile
within one condition, not a pair of trials from different conditions.
Both conditions used a common finite menu and structured final output.
Appendix~\ref{app:common-menu-protocol} identifies the selected histories,
execution contract, resource limits and recorded usage.
The optional condition permitted an observational record; the mandatory
condition required a record and let an external deterministic executor select
the action. Mandatory record--action agreement is consequently a construction
property, unlike a matching optional observational record. The combined
intervention does not isolate a pure reporting effect.

All 32 trials were evaluated: no unstarted slots, technical missingness or
protocol-invalid trials remain. This does not mean the analysis interface was
error-free: 48 attempts yielded 41 successes and seven code errors in seven
trials (five optional, two mandatory). Five errors used attribute access on
the dictionary history; two used the unavailable built-in \texttt{next}.
The prompt described \texttt{history.states}/\texttt{history.actions}, although
the worker passed a dictionary, and did not enumerate the built-ins.
Each affected trial consumed all three analysis slots, one on the error and
two successfully. Appendix~\ref{app:interface-audit} gives the trial-level
accounting with all original trials and labels retained. The effect on
subsequent actions and record structure is unassessed.
Optional P/L/U counts were 13/3/0; mandatory
counts were 10/5/1. Four purposefully chosen histories do not form a population
sample of tasks or support a model ranking. Table~\ref{tab:label-contracts}
separates the families' label contracts. The study inventory in
Appendix~\ref{app:inventory} identifies excluded historical and development
work; the later acquisition experiment has its own denominator.

\begin{table}[htbp]
\centering\small
\begin{tabular}{llrrrl}\toprule
Condition & Family & P & L & U & Verification contract\\\midrule
Optional & Continuous & 5 & 3 & 0 & Numerical, conditional\\
Mandatory & Continuous & 2 & 5 & 1 & Numerical, conditional\\
Optional & Route & 8 & 0 & 0 & Finite integer recurrence\\
Mandatory & Route & 8 & 0 & 0 & Finite integer recurrence\\\bottomrule
\end{tabular}
\caption{Native outcomes in the common-menu panel, eight trials per row.
Continuous P uses an admissible continuation witness; continuous L depends
on the numerical and historical-runtime premises in Appendix F. U remains
unresolved. Successful present-host replay does not close those premises.}
\label{tab:label-contracts}
\end{table}

\section{Results: what agreement and crossing identify}
\subsection{Pilot counterexample: native agreement does not identify the selector}
\label{sec:selector-sensitivity}
The missing choice can be seen in pilot context B. Its two archived model
analyses supplied different candidate dynamics and planning ingredients;
the native-loss representation A1 has eight candidates, while the
native-preservation representation A2 has four. The researcher reconstruction
computes one multi-step plan per candidate. To execute this component on
either representation, an analyst must also specify \emph{which plan supplies
the next action}. The declared completion R2* chooses the plan with the
smallest total penalized cost. This last selection is an analyst choice
(Appendix~\ref{app:components}), not identified by reproducing the native action.

Here \emph{native-action-compatible} means only that the selected first action
on A2 is within $10^{-6}$ of the recorded native action. It does not mean
agreement with every semantic commitment of the model's narrative or recovery
of its internal policy. We examine five selectors over the same fixed plans:
total cost, maximum target-step norm, mean squared target-step norm, trajectory
peak norm and final-step norm. These are explicit summaries of the plans'
objective, target attainment or path behavior, making the analyst's choice
inspectable. They obey the same action bounds and lowest-index tie breaking,
but need not encode the same preferences; final-step norm alone, for example,
is not the full multi-step task criterion. This is a finite retrospective
sensitivity family, not an exhaustive set of faithful completions.

\begin{table}[htbp]
\centering\small
\begin{tabular}{lccr}\toprule
Criterion minimized & Native A2: plan/label & Crossed A1: plan/label & Action norm\\\midrule
Total penalized cost (R2*) & 2/P & 2/P & 0.575465\\
Maximum target-step norm & 2/P & 2/P & 0.575465\\
Mean squared target-step norm & 2/P & 2/P & 0.575465\\
Trajectory peak norm & 2/P & 3/L & 0.042684\\
Final-step norm & 2/P & 5/P & 0.579114\\
\bottomrule\end{tabular}
\caption{Pilot context B: five selectors applied to the same fixed plan pools.
Indices are zero-based. All five select native A2 plan 2 and reproduce the
original action within $10^{-6}$. The first three select the same crossed
A1 plan; the minimum-peak selector instead selects plan 3, giving L rather
than P. Norms refer to crossed actions. Labels retain the conditional
continuous evaluation premises. These five selectors are not independent
replications or a probability sample of completions.}
\label{tab:b-completions}
\end{table}


All five select A2 plan 2, whose first action matches the native action to
approximately $1.81\times10^{-9}$. Plan 2 has a strictly lower value than
each of the other three plans on every one of the five criteria. A2 thus
provides no opportunity to discriminate among these selectors: agreement
there follows from a common best plan. Now hold the crossed A1 representation,
history and computed plans fixed. Total cost selects A1 plan 2, with action
norm $0.575465$ and a preservation witness. Trajectory peak instead selects
A1 plan 3, with norm $0.042684$ and a conditional loss certificate
(Table~\ref{tab:b-completions}). The reader command recomputes all twelve plans
and five selections, re-evaluates their first actions with the bundled grader,
and compares the complete archived witness/certificate records
(Section~\ref{sec:reader-access}). It therefore checks the label change as well
as the selector arithmetic. The numerical and historical-runtime premises
in Appendix~\ref{app:runtime-proof} remain open.

The resulting component diagnosis depends on which native-compatible
selector the analyst supplies. To use native-action agreement as a validation
test, report the candidate selectors and whether the validation cases
actually separate their choices. An independently specified selection
contract can restrict the alternatives. In this case, the common best plan
on A2 leaves that choice open. The two A1 plans also differ in action
magnitude, so this contrast evaluates the full selection change rather than
isolating a horizon or direction effect.

Under the originally declared completions, pilot context B had two P
off-diagonals and context D two L off-diagonals. Each native pair already
contained one P and one L. Those patterns describe completion-conditional
component dependence, not new outcome categories or internal causal ground
truth. The pilot's heterogeneous optimizers differ from the later common-menu
contract; the latter's structural identity is not extended to this setting.

\subsection{Common-menu applicability and recoverability}
The mandatory condition made all eight pairs executable. It did not make
their crossing tables informative by construction. Every pair used the same
declared rule, so all 32 observed-rule cells were already determined by native
actions and rule identity. Among the 16 mandatory records, 13 had one row,
one contained duplicate rows alongside distinct rows, and 13 had only one
distinct row. A duplicate-containing table can also contain distinct rows, so the
duplicate and single-distinct-row counts measure different properties.

Native actions plus the single-distinct-row check determine 29 of the
32 supported-rule-family cells. The remaining sensitivity can be calculated on each record independently.
Appendix~\ref{app:structural-control} states the fixed-contract identity and
the complete recoverability table. These calculations measure recoverable
cells; improved diagnostic accuracy requires an external diagnostic target.


Fifteen of 16 mandatory records selected the same action under both supported
rules. The remaining record changed action and its conditional continuous
P/L category; the independent
scan recovered this change. One pair involved U and remains unresolved rather
than being converted into a resolved loss. The two 32-cell displays overlap;
after within-pair deduplication they reference 30 unique numerical computations,
not 64 new samples. Cross-action labels reuse the existing bound menu census.
No new live model counterfactuals were collected for this retrospective analysis.

\paragraph{Worked same-record contrast.}
For the mandatory \texttt{probe-01} record from gpt-6-astra, repeat 1,
the independent scan yields:
\begin{center}
\begin{tabular}{lrl}\toprule
Rule applied to this record & Menu index & Bound task label\\\midrule
Uniform mean & 7 & L (conditional continuous certificate)\\
Worst case & 12 & P (continuation witness)\\\bottomrule
\end{tabular}
\end{center}
Indices are zero-based. The recorded table, menu identity and label census
are held fixed. No partner record is needed to expose this sensitivity;
it concerns the declared external rule, not a demonstrated internal cause.
The alternative action's label is a deterministic menu evaluation, not an
additional live model decision. This is the sole rule-sensitive mandatory
record, so it must not be presented as a prevalent population effect.
This trial was also affected by the dictionary-access error.
The archived table's rule sensitivity is reproducible; whether correcting
the interface would have changed that table or the decision is unknown.

The optional condition contains one submitted record, which is eligible,
and 15 trials with no submitted record (\texttt{no\_record}). It has no
eligible pairs. The fifteen absences and one eligible submission have different statuses.
With no eligible pair, the optional crossing effect remains unassessed.
The panel thus locates an applicability bottleneck in the optional condition
and a redundancy bottleneck under mandatory recording with external selection.

\section{Follow-up: required records with and without external execution}
\label{sec:followup}
The earlier mandatory condition combined a record requirement with an external
executor. The follow-up adds a required-record condition that executes the
model's proposal. Each matched triplet uses one history, repeat and
model/max/Fast profile: A permits an optional record and executes the proposal;
B requires a valid record and executes the proposal; C requires a valid record
and executes its declared selector, retaining the proposal for observation.
Every trial starts a fresh context. Four purposefully selected historical
contexts, two profiles and two repeats give 48 planned assignments, eight
triplets per profile. These repeats are not independent draws of new worlds.
The allocation and six possible within-triplet orders were fixed before
collection; earlier-panel trials are not reused as controls.
Appendix~\ref{app:followup} gives the execution segments, interface and reader
boundaries. No new population sample or live crossed decision was collected.

\begin{table}[htbp]
\centering\small
\begin{tabular}{llrrrrrrr}\toprule
Profile & Arm & Done & Tech. & P/L/U & Valid & Match & One row & Same action\\\midrule
Astra & A & 7 & 1 & 6/1/0 & 0 & -- & 0 & 0\\
Astra & B & 8 & 0 & 8/0/0 & 8 & 8/8 & 8 & 8\\
Astra & C & 7 & 1 & 5/2/0 & 7 & 7/7 & 6 & 6\\\midrule
Sol & A & 8 & 0 & 7/1/0 & 0 & -- & 0 & 0\\
Sol & B & 8 & 0 & 6/2/0 & 8 & 8/8 & 6 & 8\\
Sol & C & 8 & 0 & 7/1/0 & 8 & 8/8 & 7 & 7\\\bottomrule
\end{tabular}
\caption{Follow-up, eight planned assignments per row; profiles are separate
strata, not ranked. Done means a completed trial, not full task success.
Valid counts usable records; Match is record-selected action versus model
proposal among valid records. The empty A denominator is undefined.
Same action counts records for which both supported rules select the same
action. Sol C has one duplicate-containing record; all other rows have zero.
There are no protocol losses, active or unstarted assignments.}
\label{tab:followup-records}
\end{table}

\paragraph{Availability and correspondence.}
In A, no valid record was available among Astra's seven or Sol's eight
completed trials. This is absence of the specified external record, not
absence of reasoning. In B, each profile supplied eight valid records and all
eight matched its own proposal. Unlike C's executed-action agreement, this
correspondence was not imposed by the executor. It does not establish that
the record caused the proposal or faithfully represents internal reasoning.
Table~\ref{tab:B-structure} separates the observed record structures from
their mathematical consequences. Astra supplied eight single-row B records;
Sol supplied six single-row and two multirow B records. For every single-row
record, equality of mean and worst case is forced, so it is not an independent
empirical success of rule agreement. Sol's multirow records contain three and
four distinct rows; both also select the same action under the two rules.
That equality is an observation about these values, not forced by having one
distinct row. Their source orders are 44 and 47, respectively. No multirow
conclusion is available for Astra's empty denominator.

% Generated from the preserved minimal finals by contribution_checks.py.
\begin{table}[htbp]
\centering\small
\begin{tabular}{lrrrr}\toprule
Profile & Proposal matches & One row & Multiple rows & Multirow same action\\\midrule
Astra & 8/8 & 8 & 0 & --\\
Sol & 8/8 & 6 & 2 & 2/2\\
\bottomrule\end{tabular}
\caption{B records, separated by profile. One-row mean/worst-case equality is
forced by the contract. For multirow records, equality is an observed property
of the archived values. Astra has no multirow denominator, so its dash is
undefined, not a success. Sol's two multirow records have three and four
distinct rows, respectively.}
\label{tab:B-structure}
\end{table}


The empirical result is the co-occurrence of proposal matching and these
record structures without external selection. The action-copy control in
Section~\ref{sec:measurement} establishes why matching alone leaves
substantive use of the history unresolved. The actual records' predictive
value requires its own measurement. Single-row and duplicate records satisfy
the protocol and remain in the structural accounting. For evaluation practice, a correspondence rate should
therefore be accompanied by the number of distinct rows, the actual choices
under alternative supported rules, and the action-copy control. These checks
separate an observed match from an opportunity to discriminate between rules;
they do not replace a test of prediction on unseen decisions.

\paragraph{Conditional outcomes and observed opportunities.}
Table~\ref{tab:followup-pairs} reports the two predeclared contrasts using only
observed nontechnical pairs. Astra has seven fully completed triplets and Sol
eight. L outcomes remain in the comparison; technical failures are reported
separately and are not converted into L. The directions do not support a
common improvement claim across these profiles, nor an equivalence or
no-effect conclusion. In Sol B, eight matching records coexist with two L
labels, illustrating why correspondence and continuation feasibility require
separate measurements. P certifies a conditional feasible continuation;
the experiments do not measure completed-task success.

% Generated from identity-bound existing follow-up sources.
\begin{table}[htbp]
\centering\small
\begin{tabular}{llrrrrrr}\toprule
Profile & Contrast & Pairs & Missing & L$\to$P & P$\to$L & Same & Net P\\\midrule
Astra & B--A & 7/8 & 1 & 1 & 0 & 6 & +1\\
Astra & C--B & 7/8 & 1 & 0 & 2 & 5 & -2\\
Sol & B--A & 8/8 & 0 & 0 & 1 & 7 & -1\\
Sol & C--B & 8/8 & 0 & 2 & 1 & 5 & +1\\
\bottomrule
\end{tabular}
\caption{Descriptive paired transitions run from A to B for B--A, and from B to C for C--B. Pairs are observed/planned; Missing counts technical missing pairs; Same includes P-to-P and L-to-L. Dividing Net P by observed pairs recovers the reported paired difference. No pooled estimate or population inference. C--B is the total condition contrast, not an isolated executor effect on the same fixed record.}
\label{tab:followup-pairs}
\end{table}


\paragraph{Where the outcome contrasts occur.}
The family breakdown in Table~\ref{tab:followup-families} locates every
observed P/L movement in the two continuous-control histories. All completed
route decisions were P, although each route menu contained an L option;
this pattern was not forced by having only P choices. The continuous menus
have 13 actions and the route menus three, so these purposeful contexts are
not a calibrated difficulty ladder. Appendix~\ref{app:followup-history}
reports every profile/history/arm cell. In particular, Sol's C--B net gain
of one P comprises two L-to-P and one P-to-L transitions. Net change alone
does not describe those opposing movements, and the underlying continuous
labels retain the numerical/runtime premises stated above.

% Generated from identity-bound existing follow-up sources.
\begin{table}[htbp]
\centering\small
\begin{tabular}{llrrrr}\toprule
Profile & Family & Assigned & Done & Tech. & P/L/U\\\midrule
Astra & Continuous & 12 & 12 & 0 & 9/3/0\\
Astra & Route & 12 & 10 & 2 & 10/0/0\\
Sol & Continuous & 12 & 12 & 0 & 8/4/0\\
Sol & Route & 12 & 12 & 0 & 12/0/0\\
\bottomrule
\end{tabular}
\caption{Follow-up outcomes by family, separately within each profile. Each family has two purposeful histories and two repeats per arm. Done denotes a completed decision, not full task success. P/L labels in the continuous family retain the numerical/runtime premises.}
\label{tab:followup-families}
\end{table}


\paragraph{A changed action need not change the task label.}
The follow-up contains two rule-sensitive records, both in C
(Table~\ref{tab:followup-rule-examples}). Mean and worst case select different
menu indices, while the bound label remains P in one and L in the other.
These are existing-record calculations, not new live counterfactual trials.
Together with the older panel's single label-changing example, they make
action sensitivity and task-label sensitivity separately inspectable.

% Generated from identity-bound existing follow-up sources.
\begin{table}[htbp]
\centering\small
\begin{tabular}{llrlll}\toprule
Profile & History & Repeat & Arm & Mean & Worst case\\\midrule
Astra & probe-01 & 2 & C & 1 (P) & 0 (P)\\
Sol & probe-02 & 1 & C & 11 (L) & 0 (L)\\
\bottomrule
\end{tabular}
\caption{The two rule-sensitive follow-up records. Entries give zero-based menu index and bound label. The action changes but the label does not; these are per-record calculations on existing menus, not live crossed decisions. Continuous labels remain conditional.}
\label{tab:followup-rule-examples}
\end{table}


Every completed C proposal already matched the record-selected action
(Astra 7/7, Sol 8/8). No observed C trial therefore tests the benefit of
overriding a conflicting proposal. The C--B contrast is not an isolated
mechanical correction effect on the same fixed record. Labels and
within-record rule diagnostics are tied to the existing menu contract;
they are not additional live crossing observations or independent evidence
about new task physics.

\section{Prospective pilot: does a frozen record predict a new choice?}
\label{sec:prospective}
We test whether a source decision record improves prediction of later model
choices beyond the public history and source action. Four new purposefully
constructed histories in the two disclosed task families are used separately
for Astra/max/Fast and Sol/max/Fast. Each source has two precommitted target
branches: eight target choices per profile, grouped in four histories.
The histories, records and target choices form a new panel. Each target
extends its source history by one observed transition under a fixed exogenous
action: continuous sources at step 6 branch to step 7 under $(0.6,0)$ or
$(-0.6,0)$; route sources at step 4 branch to step 5 by advancing one or three
nodes. These alternatives were set before source-model outputs and need not
follow the source proposal. Table~\ref{tab:prospective-branches} and
Appendix~\ref{app:prospective} give the construction and selection basis.

\paragraph{Information and output protocols.}
The source actor's final response contains an action, a numeric table with a
declared rule, and four brief fields for assumptions, uncertainty, criterion
and expected observation. The record is frozen before the source actor sees
either target. Two fresh-context observers of the same requested profile
predict the target actor's top-1 action with or without that record. Both
receive the source action, source and target histories, menus, and the
complete source and target instructions and output schemas. Only the record
differs between observer inputs. Their order is preallocated, allowances are
equal, and both predictions are locked before the fresh target actor runs.

The target's final response contains an action alone. Source and target may
both request up to three local analyses before their final response, and
their proposals are unforced. The target receives the public evidence and
source action, without the record or observer predictions. The comparison
therefore measures the record's added value for predicting an action under
this target protocol. Both observers predict the same target; the effect of
changing the source--target output protocol is not separately identified.

\begin{table}[htbp]
\centering\small
\begin{tabular}{llrrrrrl}\toprule
Panel & Profile & Pairs & Record & No record & Copy & Changed & Difference\\\midrule
Initial & Astra & 8/8 & 6/8 & 7/8 & 5/8 & 3/8 & -0.125\\
Initial & Sol & 7/8 & 3/7 & 3/7 & 0/7 & 7/7 & $[-0.125,+0.125]$\\
Combined & Astra & 8/8 & 6/8 & 7/8 & 5/8 & 3/8 & -0.125\\
Combined & Sol & 8/8 & 3/8 & 3/8 & 0/8 & 8/8 & 0.000\\
\bottomrule\end{tabular}
\caption{Prospective top-1 action prediction. Pairs retains the eight assigned
targets/profile; other fractions use available valid targets. Difference is
the equally weighted mean over four history clusters, with two target slots
each. The initial Sol interval is a missing-data identification bound over all
eight assigned targets, not a confidence interval; its observed 3/7 versus
3/7 comparison has difference zero. Copy uses the old source action; Changed
compares the target action with that action. Profiles are not ranked.}
\label{tab:prospective}
\end{table}

\begin{table}[htbp]
\centering\small
\begin{tabular}{lrrrr}\toprule
Profile & Both correct & Record only & No-record only & Both wrong\\\midrule
Astra & 5 & 1 & 2 & 0\\
Sol & 2 & 1 & 1 & 4\\
\bottomrule\end{tabular}
\caption{Joint top-1 prediction correctness for the combined panel, eight
targets per profile. Record only and No-record only identify which observer
is correct on the same target. Each cell uses one generation per role;
the counts describe paired outcomes, not case-level causal effects.}
\label{tab:prospective-paired}
\end{table}


\paragraph{Prediction results.}
The combined panel yields 6/8 versus 7/8 correct predictions for Astra and
3/8 versus 3/8 for Sol, with versus without the record. The mean of the four
within-history paired differences is $-0.125$ for Astra and $0$ for Sol.
Table~\ref{tab:prospective-paired} shows the joint outcomes: Sol's equal
totals include one target correctly predicted only with the record and
another correctly predicted only without it. Across these assignments,
no positive net accuracy gain is observed. Each cell uses one generation
per observer and target role; a discordant pair does not by itself identify
a case-level causal benefit or harm.

\paragraph{Where the prediction contrasts occur.}
Table~\ref{tab:prospective-families} adds a post hoc descriptive breakdown.
Sol's aggregate zero combines one additional correct prediction with the
record on continuous tasks and one fewer on route tasks. Astra's one fewer
correct prediction occurs in continuous control; both observers predict all
four route targets correctly. Each family contains two purposeful histories
with two linked targets each. The prespecified per-profile aggregate remains
primary; these counts locate the observed contrasts without estimating a
family-level causal effect. Table~\ref{tab:prospective-histories} gives all
eight history/profile units.

\begin{table}[htbp]
\centering\small
\begin{tabular}{llrrr}\toprule
Profile & Family & With record & Without record & Net correct\\\midrule
Astra & Continuous & 2/4 & 3/4 & $-1$\\
Astra & Route & 4/4 & 4/4 & $0$\\
Sol & Continuous & 3/4 & 2/4 & $+1$\\
Sol & Route & 0/4 & 1/4 & $-1$\\
\bottomrule\end{tabular}
\caption{Post hoc descriptive family breakdown of the combined prospective
panel. Each family has two purposeful histories with two linked targets each.
Net correct is the count with minus without the record. The prespecified
per-profile aggregate remains primary; these four-target contrasts do not
estimate family-level causal effects.}
\label{tab:prospective-families}
\end{table}


Copying the source action predicts 5/8 Astra and 0/8 Sol targets. The target
changes action relative to the source in 3/8 and 8/8 cases, respectively.
This baseline describes source--target action change; variability across
repeated runs of an identical target input was not measured.

\paragraph{Record structure and interpretation.}
All four source records per profile match the unforced source proposal.
Seven of the eight numeric tables contain a single row: four for Astra and
three for Sol. Agreement of the two supported rules is forced in these seven
tables and observed in the remaining two-distinct-row Sol table. No duplicate
row occurs. The semantic informativeness of the accompanying prose was not
independently evaluated. The treatment is the complete record, including its
prose; table row counts alone cannot explain the prediction result.

The reference controls verify protocol handling and prediction of supplied
hand-authored policies (Appendix~\ref{app:prospective}). Sensitivity to an
informative record in the scientific tasks remains a separate control target.
The observed contrast combines record information, observer use, the
source--target protocol change and one draw of each observer and target.
It does not separate those explanations for the net result. Task-matched
informative-record controls, a same-output-protocol comparison and repeated
identical-input decisions would distinguish them in a further study.
With four purposeful histories per profile and two linked targets per
history, the present pilot characterizes these assignments.

\paragraph{Execution segments.}
The original run produced fifteen targets. After inspecting its partial
results, we completed the remaining assigned Sol block in a separate segment
with the same source, inputs, profile and allowances and four fresh reference
checks. Table~\ref{tab:prospective} retains both panels. The completion adds
one target to an existing history; no completed scientific decision was
rerun. Appendix~\ref{app:prospective} retains the interrupted attempt, all
assigned roles and usage. The two segments may differ in time or provider
state; they form one descriptive pilot, not a confirmatory replication.

\section{Position relative to existing work}
\label{sec:related}
\paragraph{Execution and component diagnosis.}
\citet{lyu2023faithful} translate a problem into an intermediate representation
and use a deterministic solver to obtain an answer, locating their faithfulness
guarantee in the solution stage. Our C condition similarly binds execution
to an external record. We examine the additional evaluation questions raised
by constructing and comparing such records: selector-completion sensitivity,
correspondence without enforced execution, and prediction of later choices.
CAR estimates effects of interventions on agent steps through same-policy
downstream replay \citep{shah2026car}; AgentRx uses constraint evidence and a
judge for execution diagnosis \citep{barke2026agentrx}; and BLADE aligns evidence
from action stages for inspection \citep{guo2026blade}. RUU's component
diagnoses instead attach feasibility labels to actions selected from
external records. The diagnosed object and the outcome measure therefore
differ from stochastic-policy replay or failure attribution.

\paragraph{Alternative explanations and discriminating tests.}
\citet{jain2019attention} show that substantially different attention
distributions can preserve model predictions. \citet{wiegreffe2019attention}
argue that explanation claims depend on the model relationship being tested,
and introduce baselines, variance calibration and diagnostic tests to assess
attention's usefulness. Their debate distinguishes finding an alternative
from evaluating the explanatory value of the original. Our counterexample
concerns an analyst's external selector completion: the common best plan
on A2 makes the original-action match nondiscriminating, while A1 separates
the completed selectors. It identifies an ambiguity in that validation
criterion, rather than testing attention or establishing that the original
record is uninformative. This scope follows the broader requirement to
evaluate explanations against the system they describe
\citep{jacovi2020faithful}. Interventions on neural representations address
internal causal relations directly \citep{geiger2021causal}; our intervention
is on the external record and selection rule.

\paragraph{Prediction-based explanation evaluation.}
Leakage-Adjusted Simulatability (LAS) measures whether an explanation helps an
observer predict a model output while controlling for direct output leakage
\citep{hase2020las}. REV measures new label-relevant information beyond an
input-and-label baseline using conditional V-information \citep{chen2023rev}.
These approaches motivate assessing what an explanation adds to information
already available. \citet{madsen2024self} evaluate self-explanation consistency
through counterfactual, attribution and redaction tasks, while
\citet{mayne2026positive} test added explanation value for predicting a
reference model's answers on related inputs. They also use repeated
reference-model outputs to assess the prediction ceiling imposed by target
variability (their Section 4.3 and Appendix A.7); our single-generation pilot
does not estimate that ceiling.

Our prospective pilot applies the added-prediction question to records frozen
before precommitted target branches. Both observers already receive the source
action; neither sees the target choice before predicting it. The measured
quantity is the paired top-1 accuracy difference on later choices across
different source and target output protocols. LAS's leakage adjustment and
REV's information metric are distinct estimands. The pilot supplies a
source-auditable RUU case for this established evaluation question, alongside
record-structure and source-action-copy controls.

\paragraph{Connection to scientific discovery.}
BoxingGym evaluates experiment selection through expected information gain
and model discovery through predictions, including whether a scientist's
explanation helps another agent predict without further experiments
\citep{gandhi2025boxinggym}. RUU's present study concerns interpretation of
decision records and prediction of model choices. Prediction of withheld
environment observations and the utility of selected experiments are the
next behavioral targets in our broader research program. Keeping these
outcomes distinct connects the measurement lesson here to discovery
evaluation without treating action correspondence as discovery success.

\section{Interpretation and next research questions}
The case study yields a practical sequence for evaluating decision records.
First, establish which records support the proposed operation. The transfer
panel's zero eligible pairs locate an applicability bottleneck. Second,
specify the analyst's completion and test whether the validation cases
separate plausible alternatives. The pilot's common best plan shows why a
successful original-action match can leave this choice open. Third, report
correspondence together with structural opportunities for rule discrimination.
The unforced-proposal follow-up demonstrates their coexistence: all sixteen
required records match, while fourteen have mathematically forced rule
agreement. Finally, evaluate added prediction on new decisions. The prospective
pilot supplies that separate measurement and observes no positive net accuracy
gain in its assigned cases.

These checks also determine the appropriate unit of evidence. Under the fixed
numeric contract, a per-record rule scan exposes the sensitivity in the paired
display (Appendix~\ref{app:structural-control}). Action changes and P/L/U label
changes answer different questions, as the two rule-sensitive but
label-preserving records illustrate. The proposed audit makes these
distinctions inspectable; its effect on downstream diagnostic accuracy remains
to be evaluated against independently established diagnoses.

\paragraph{Scope of the evidence.}
The histories are purposefully selected within two task families. Continuous
P/L labels retain the numerical and historical-runtime premises in
Appendix~\ref{app:runtime-proof}; P certifies a feasible continuation, not a
completed task. All completed C proposals already agree with their selectors,
so correcting a conflicting proposal has no observed denominator. The older
panel's seven analysis errors, including its sole label-changing
rule-sensitive record, have no isolated behavioral-effect estimate. Record
correspondence and external substitutions address observed outputs; internal
causal faithfulness requires a different intervention. In the prospective
pilot, same-protocol targets, repeated identical-input targets, and a control
that demonstrates use of informative records in the scientific tasks would
help distinguish protocol transfer, target variability and observer sensitivity.

\paragraph{Research under uncertainty.}
Three questions guide the next stage: how a stated assumption relates to the
next experiment choice; whether that experiment resolves more uncertainty
than an equally resourced valid fixed plan; and whether a revised explanation
improves predictions for cases withheld during its construction. Each question
calls for its own scored behavioral outcome. The present prospective pilot
measures model-choice prediction. Prediction of new world observations and
the utility of adaptive experimentation will extend this program from
interpreting decision records to evaluating scientific behavior.

\section{Reader access and reproducibility}
\label{sec:reader-access}
The anonymous supplement supplies the manuscript sources, current-study
mechanisms, visible interaction evidence and source-bound readers.
Appendix~\ref{app:reader-access} lists the material and check boundaries for
each evidence set. One archive contains both the main and prospective
components in a lossless text representation; its reader restores their
exact bytes, verifies their manifests and runs the evidence checks.
From the extracted supplementary archive, run:
\begin{verbatim}
python3.12 -I -B -S verify.py
\end{verbatim}
The command recomputes the reported source-level results, checks the paired
prediction outcomes and selector-discrimination example, and runs the
existing numerical replays. No model account is needed. The README also
provides permanent extraction and corruption-test commands. Source
recalculation and present-host numerical replay establish reproducibility
of the supplied evidence; historical/provider attestation and independent
laboratory replication remain separate verification levels.

\paragraph{Responsible interpretation.}
The experiments use simulated environments and model outputs. The checks
evaluate external decision records; assessing internal reasoning, intentions
or deployment safety requires evidence for those distinct claims. The
disclosed mechanisms support inspection of this fixed study. Future
benchmark mechanisms have a separate disclosure policy.

\bibliography{references}
\bibliographystyle{tmlr}
\clearpage\appendix
\section{Disclosed task, mechanism and delivered information}
\label{app:task}
Represent the two-dimensional state and action by complex numbers $z_t$ and $u_t$. A step is indexed by its pre-action state: $u_0$ produces $z_1$. The initial conditions are $z_0=1$, $\phi_0=4.147589101011308$ and $u_{-1}=0$. The evaluated world obeys
\begin{align}
 d(z)&=(1.06+0.015|z|^2)z+(0.12-0.08i),\\
 z_{t+1}&=d(z_t)+e^{i\phi_t}u_t,\qquad |u_t|\leq 0.6,\\
 \phi_{t+1}&=\phi_t+\omega+\kappa\,\operatorname{Re}(u_{t-1})\operatorname{Im}(u_{t-1}),\qquad\omega=0.1.
\end{align}
The lag in the phase update is part of the mechanism: the update uses $u_{t-1}$, not the action just chosen. The original family allowed $\kappa\in\{0.3,3,6,12,18\}$. Only the two settings below occur in the fixed-history pilot. They are parameter settings, not a validated difficulty ordering.

\begin{center}
\begin{tabular}{llll}
\toprule
Paper alias & Archival case identifier & Snapshot after step & $\kappa$\\
\midrule
A & case-01 & 4 & 12\\
B & boundary-01 & 11 & 12\\
C & case-03 & 4 & 18\\
D & boundary-02 & 9 & 18\\
\bottomrule
\end{tabular}
\end{center}

An episode ends on escape ($|z_t|>4$) or after 32 actions. Success requires completion and $|z_t|\leq0.05+10^{-10}$ at \emph{all} $t\in\{28,29,30,31,32\}$. Merely avoiding escape is insufficient. Action validation rejects nonfinite coordinates or norms above $0.6+10^{-9}$; a smaller tolerance overshoot is normalized to $0.6$. State observations are rounded to twelve decimal places, whereas the simulator retains binary floating-point state. Stored actions retain their applied precision. The exact numeric history of every selected case is included in \texttt{legacy-replay/data/cases.json}; it is not regenerated from a caption or rounded manuscript table.

The participant receives the objective, horizon, action limit, escape threshold, full permitted action--state history and remaining analysis allowance. It does not receive these equations, $\phi_t$, $\omega$, $\kappa$, the evaluator's continuation or its loss diagnosis. The opaque identifiers in delivered inputs do not disclose the boundary/comparison labels. The frozen selector and schemas are inspectable in the reader artifact; the full provider logs are not redistributed. Disclosure to reviewers does not alter the historical participant information boundary.

\section{Witness and loss-bound justification}
\label{app:bounds}
\paragraph{Positive witness.}
For every pre-action snapshot and every resulting state, one fixed privileged controller attempts a continuation. At each step it chooses
\begin{equation}
 u=\operatorname{clip}_{0.6}\!\left[-d(z)e^{-i\phi}\right],\qquad
 \operatorname{clip}_{r}(v)=
 \begin{cases}
   v, & \lVert v\rVert\leq r,\\[3pt]
   \displaystyle\frac{r v}{\lVert v\rVert}, & \lVert v\rVert>r.
 \end{cases}
\end{equation}
Here $\lVert v\rVert$ is the Euclidean norm, equal to the modulus for a complex scalar.
The original transition, action validation, escape check and scheduled target test score the entire continuation. A successful admissible continuation establishes existence for the implemented task. A failed continuation does not establish impossibility. This controller knows the hidden phase and mechanism; its success is not evidence of blind action identification or participant learning. Every pilot pre-action state passed this positive baseline.

\paragraph{Sufficient negative bound for the ideal recurrence.}
Write \(x = (\operatorname{Re} z, \operatorname{Im} z)\) and $c=(0.12,-0.08)$. Fix a vector $q$ with $\|q\|_2\leq1$. Suppose $q^Tx\geq p>0$. Since $\|x\|_2\geq q^Tx\geq p$, and the scalar multiplier $1.06+0.015\|x\|_2^2$ is positive, the next-state projection satisfies
\begin{align}
q^Tx'&=(1.06+0.015\|x\|_2^2)q^Tx+q^Tc+q^TR(\phi)u\\
&\geq (1.06+0.015p^2)p+q^Tc-0.6.
\end{align}
The last term follows from Cauchy--Schwarz and the admissible action radius. The function $g(p)=(1.06+0.015p^2)p$ is increasing on the positive domain. Consequently, iterating this lower bound while it remains positive gives a bound for \emph{every} allowed future control sequence, irrespective of the future phases. A bound exceeding 4 before the horizon implies unavoidable escape; a bound exceeding the target tolerance at a required target step excludes successful completion. A target violation already present in the recorded prefix also suffices, because future actions cannot change a past required state.

\paragraph{Implementation and numerical scope.}
The evaluator chooses $q$ by truncating each coordinate of $x/\|x\|$ toward zero to four decimal places, and checks its norm. It initializes $p_0=q^Tx$ from the stored full-precision state. The implemented recurrence is
\begin{equation}
p_{j+1}=\left\lfloor (1.06+0.015p_j^2)p_j+q^Tc-0.6-10^{-9}\right\rfloor_{10^{-12}},
\end{equation}
where $\lfloor\cdot\rfloor_{10^{-12}}$ rounds downward to that grid. Coefficients and initial state are converted from their actual binary-float values to exact rational numbers. Loss is reported only for a positive induction sequence crossing $4.000001$ within the remaining horizon, or crossing $0.05+10^{-10}$ at a required target step. A separate 110-digit decimal implementation checks the direction, every bound, the positivity domain, threshold and horizon. The code also checks that the transition definition and task constants match the bound's assumptions.

\paragraph{Archived arithmetic audit.}
A subsequent record-only audit covers the pilot, distinct component
grades and the continuous subset of the older 48-trial transfer panel. All 50 distinct
stored directions, including directions for which no loss was certified,
have exact squared norm at most one; the minimum value of
$1-\|q\|_2^2$ is $5.75\times10^{-6}$. There are ten distinct certified
projection records (eleven occurrences). Rational recomputation reproduces
the archived binary64 bounds exactly. Recomputing these recurrences with per-step
subtractions $10^{-9}$, $2\times10^{-9}$ and $10^{-8}$ preserves every
certificate kind and first certification step. Even at the enlarged
slacks, the smallest final lower-bound gap above the applicable runtime
threshold exceeds $0.0371$. No simulator transition or model request is
performed by this audit, and no historical certificate is rewritten.

The action validator rejects computed norms above $0.6+10^{-9}$ and
rescales a smaller computed overshoot by $a\,0.6/|a|$. It therefore does
not simply apply every action up to the rejection tolerance unchanged.
For the 1,201 distinct applied actions in the inspected continuous
trajectory records, exact squared-norm checks give a maximum norm of
approximately $0.60000000000000002694$. This finite record check does not
bound all possible future actions or transcendental-function errors.

\paragraph{Runtime-specific conditional budget.}
A separate local audit inspects the pinned interpreter and the currently
loaded ARM64 math image. Its compiled complex multiplication uses fused
multiply-add instructions; its complex exponential calls
\texttt{exp} and \texttt{\_\_sincos\_stret}. The norm uses platform
\texttt{hypot}, and the square still uses platform \texttt{pow}.
Thus a source-level list of separate sine/cosine calls does not fully
describe this executable.

Appendix~\ref{app:runtime-proof} derives a norm bound for the captured
\texttt{hypot} instruction sequence under ordinary IEEE binary64
round-to-nearest semantics with gradual underflow. It then reduces the
remaining library accuracy requirements to an absolute square error
$\epsilon_{\rm pow}$ and a norm excess $\delta_{\rm rot}$ for the
computed complex rotation. Accuracy relative to the true angle is not
needed for this projection bound. Taking each of these two bounds as
$2^{-30}$ yields total one-step error below $6.288\times10^{-10}$,
which fits the archived $10^{-9}$ subtraction. The local checker verifies
the scalar inequalities with exact rational arithmetic.

This is a sufficient conditional bound, not an attestation that the two
library bounds hold uniformly. Finite current-process probes satisfy them
but do not prove them. In addition, the two inspected historical runtime
inventories and activation records do not pin the loaded math-image
identity or floating-point control register. Today's image capture cannot
retroactively fill that gap. The ten archived projection certificates
therefore retain their stated conditional numerical status. The route
family's exact integer transitions and finite-horizon dynamic program do
not use this projection premise.

The classifier returns preserved only with a successful witness, lost only with a sufficient negative certificate, and unresolved otherwise. Simultaneous positive and negative evidence is treated as an error. This is an incomplete classifier, not a complete viability kernel or a map of a continuous safe cone. A numerical failure in candidate optimization is a different unresolved category and is not fed to this classifier as physical loss.

\section{Computational interventions and reproducible selectors}
\label{app:components}
The candidate coefficients and original analysis outputs are fixed data. No candidate is refitted after a component is crossed. Let $f_m+B_m u$ denote candidate $m$'s predicted next state at the selected snapshot. Candidate construction, feature families, phase extrapolation and the exact numeric coefficients are inspectable in \texttt{legacy-replay/data/component\_inputs.json}. The replay implements these fixed operations; it does not execute the archived participant programs.

\paragraph{One-step worst-case rule.}
For $G$ equally spaced angles and radius $r$, set $u_j=r(\cos(2\pi j/G),\sin(2\pi j/G))$ and choose the first grid index minimizing $\max_m\|f_m+B_m u_j\|_2$. Snapshot B uses $G = 720$ and $r = 0.6$; D uses $G = 1440$ and $r = 0.599999999$. The four finer-grid checks use $G=2880$. This is optimization on a finite circle, not a guarantee of global optimization over the full action disk or of task completion.

\paragraph{Past-error weighted rule.}
Each fitted parent has weight $w_m=1/(10^{-4}+e_m^2)$, where $e_m$ is its recorded leave-one-out residual RMS. In the expanded D component, six scenarios inherit equal shares $w_m/6$ of each parent's weight. With $v=\sum_m w_m B_m^Tf_m$, the selected action is $u=-0.6v/\|v\|$. The matrices here are scalar gains times rotations, so the quadratic term in $\sum_mw_m\|f_m+B_mu\|^2$ is constant on the radius-$0.6$ circle. Thus this direction minimizes that fixed-radius surrogate for these matrices. It is not a general theorem for arbitrary linear-control matrices, a calibrated posterior, or a worst-case task-safety guarantee. D's second representation comprises 11 parent fits expanded by two phase modes and three phase offsets $\{-0.12,\,0,\,0.12\}$; the complete construction is in the code and numeric inputs.

\paragraph{Multi-step plans and the researcher completion.}
For each of B's candidates, the individual planner optimizes the remaining $N=32-n$ controls with projection to the action disk. Its surrogate objective is
\begin{equation}
 J_m=\sum_{i=1}^{N}\left[w_{n+i}\|\hat x_{n+i}\|^2+\lambda\big(\max\{0,\|\hat x_{n+i}\|^2-14.44\}\big)^2+10^{-5}\|u_{n+i-1}\|^2\right],
\end{equation}
where $w_t=1$ for $t\geq28$ and $10^{-4}$ otherwise, and $\lambda=10$. The deterministic initialization, projected accelerated-gradient updates, backtracking limits and restart schedule are included in \texttt{components.py}; individual plans use 1,000 iterations. The declared selector R2* chooses the individual plan with the lowest computed cost and takes its first action. This final selection rule was completed by the researchers. Agreement with the original choice is a reconstruction check, not proof that it was the model's latent policy. At least one crossed action has magnitude approximately $0.5755$, so this contrast does not isolate direction or horizon while holding magnitude fixed.

The two auxiliary joint plans use a shared first action, separate candidate-specific later actions, the mean of the candidate costs, and 1,400 iterations. One is numerically unresolved. The one-at-a-time D deletions remove each of 66 scenarios or each of 11 six-scenario parent groups; all other settings stay fixed and no refitting occurs. No-rescue results cover only these interventions, not joint deletion, changed fitting or arbitrary alternative uncertainty models.

The 91 planned cells comprise 8 primary crossings, 4 finer-grid checks, 2 joint-plan auxiliaries and 77 omissions. Ninety were computed and one auxiliary remained numerically unresolved. Repeated actions are deduplicated before privileged grading, giving 13 distinct grades. There is no additional model sample in this count. The full records, rather than the rounded values in figures, define action identity.

\subsection{Five retrospective completions at pilot context B}
\label{app:b-completions}
This separate sensitivity analysis holds the eight individual plans from
B's native-loss pool A1 and four from its native-preservation pool A2 fixed.
It changes only which plan
supplies the first action. For a plan's predicted states $\hat x_t$, the
five minimized criteria are: the total penalized cost $J_m$ above;
$\max_{28\leq t\leq32}\|\hat x_t\|$;
$\tfrac15\sum_{t=28}^{32}\|\hat x_t\|^2$;
$\max_{n\leq t\leq32}\|\hat x_t\|$, including the common starting state;
and $\|\hat x_{32}\|$. All ties use the lowest existing plan index.
The last criterion covers only the final target step; it is not the full
task-success condition. The finite selector list was fixed before its
additional grading, with the original outcomes already known.

Table~\ref{tab:b-completions} reports these five selectors in the main text.


Thus native-action reproduction alone does not identify the crossed
completion. The reversal also changes the chosen plan and action magnitude;
it does not isolate a horizon-only or direction-only effect. The saved
selector indices, exact actions, metrics and labels are in
\texttt{artifact/source/pilot-sensitivity.private.json}. The original
action lock and full twelve grade records, including the fixed-controller
continuation trajectories and projection certificate, are supplied alongside
it as \texttt{pilot-sensitivity-actions-lock.json} and
\texttt{pilot-sensitivity-grades.json}. Their original hashes are checked
before replay. The standard reader uses recomputed plan actions as grader
inputs, binds the case, history, world, action and certificate state, and
compares the complete replayed grade with the archive. Discrete fields and
identities must match exactly; floating trace values and plan metrics use
$|a-b|\leq10^{-8}(1+|b|)$, while native-action compatibility uses $10^{-6}$.
Each projection certificate also passes the separate 110-digit arithmetic
check; this does not prove the uniform library bounds or historical runtime
premises in Appendix~\ref{app:runtime-proof}. The reader reports its twelve
verification grade calls separately from the original collection counters.
This finite
sensitivity result does not cover new optimizers, mixtures or arbitrary
alternative policies.

\section{Protocol, record boundaries and cost accounting}
\label{app:protocol}
The fixed trial order was A1, D1, C1, B1, C2, B2, A2, D2, where the suffix denotes repeat. Each trial begins in a fresh local context with an opaque snapshot identifier; both repeats of a case receive exactly the same initial payload bytes. Each subsequent payload contains the same physical history and that trial's own submitted analysis code and returned results. Therefore the third or fourth requests are not equal-input repetitions of the first request.

The wire schema permits one analysis request or one final two-coordinate action. Analysis code must define \texttt{act(history)} and can use the allowed mathematical facilities; no network, filesystem, process creation or hidden simulator interface is provided. The protocol and policy source state the exact permitted syntax and names. Local analysis has 30-second CPU/wall limits, a sampled 1 GiB RSS stop and a 100,000-character output limit. Sampled RSS is not a hard allocation ceiling. Each trial permits at most three analysis requests and one final action, for at most four model requests in total. Trial wall time is 3,780 seconds; request time is bounded by the smaller of 900 seconds and remaining trial time minus a 120-second record/grading margin. No retry, neutral request, replacement or outcome-dependent stopping is permitted.

The model identifier and effort are configuration records, not attested server weights. Local receipts record one agent message per call, no unexpected tool events and cleanup of initially empty call directories. These checks do not establish independent sampling seeds or provider-side isolation. Full provider transcripts remain in the source archive; this bundle distributes selected numerical inputs and final record projections. Hidden chain-of-thought is not observed or claimed as evidence.

The recorded 588,051 input and 129,677 output tokens total 717,728. The 92,160 cached-input tokens are already included in input. Reasoning-output tokens are likewise a subset of output. There were 28 receipted requests and zero unreceipted requests in this pilot. Provider receipts are the source of these counts; analyst use and money are unmeasured. The historical technical null's unresolved reservation is not a zero-use observation and is not included in the pilot total.

\section{Prospective tasks, reference histories and adapter}
\label{app:transfer}
Exact configurations and twelve packets are in
\texttt{legacy-replay/data/transfer/}; the numeric replay inputs under
\texttt{legacy-replay/data/transfer/} are checked against these records.

\paragraph{Continuous reference histories.}
The recurrence and task target are those in Appendix A, with three new
initial phase/frequency configurations and couplings 3, 6 and 12.
Phases and frequencies derive deterministically from SHA-256 of the
literal seed \texttt{transfer-validation-20261025-v1} and instance
ordinal. A prescribed privileged ten-step path tracks waypoints with
radii $1.1,1.2,1.3,1.4,1.5,1.6,1.75,1.9,2.05,2.2$ and deterministic
changing angles. Snapshots are after steps 6 and 10. Pre-collection
checks require legal histories and a positive continuation, without
resampling. This selection does not establish calibrated boundary
distances or natural participant behavior.

\paragraph{Route--resource dynamics and certificates.}
The state is node $n\in\{0,\ldots,18\}$ and remaining integer fuel $f$.
Action $j\in\{0,1,2\}$ is encoded as $[j,0]$ and advances to
$n'=\min(18,n+j+1)$. The hidden stationary positive edge cost is
\begin{equation}
c(n,j)=b_j+((n+o_j)\bmod p_j),\qquad f'=f-c(n,j).
\end{equation}
Success means reaching node 18 within twelve actions without negative
fuel. The three fixed $(b,o,p)$ configurations and initial fuel values
are in \texttt{CASES.private.json}. Unlike continuous control, every
legal route action is in the menu. The participant knows the graph,
objective, horizon and observed costs through its action--state history,
but not the edge-cost formula or private coefficients.

Let $V_h(n)$ be the minimum fuel needed to reach the goal in at most
$h$ further actions. Set $V_0(18)=0$, $V_0(n<18)=+\infty$, and
\begin{equation}
V_h(18)=0,\qquad
V_h(n)=\min_{j\in\{0,1,2\}}
\{c(n,j)+V_{h-1}(\min(18,n+j+1))\}.
\end{equation}
The code represents infinity by \texttt{None}. With positive costs,
$f\geq V_h(n)$ provides a minimum-cost admissible continuation;
$f<V_h(n)$, negative fuel or an unreachable goal certifies loss.
Every value-table cell is checked by a separately implemented recursive
recurrence. Positive witnesses are replayed. Initial fuel is the exact
minimum completion cost plus one; the lexicographically first
minimum-cost path supplies snapshots after steps 2 and 4. Baseline
feasibility does not prove that the hidden cost law is identifiable from
the observed short prefix.

\paragraph{Exact optional table contract.}
The diagnostic continuous menu has twelve equally spaced directions at
radius $0.6$, rounded to twelve decimals, plus zero. For a returned
table with losses $L_{mj}$ and weights $w_m$, the weighted rule chooses
the first menu index minimizing $\sum_m w_mL_{mj}/\sum_mw_m$; the
worst-case rule minimizes $\max_m L_{mj}$ with the same tie break.
Weights belong to the candidate pool. The adapter accepts only the
exact keys \texttt{actions}, \texttt{candidate\_losses},
\texttt{weights}, and \texttt{rule}, the exact delivered menu, 1--128
candidate rows and finite entries of absolute magnitude at most
$10^{100}$. Weights are finite, nonnegative and bounded by $10^{100}$,
with positive total. Final-action reproduction uses maximum coordinate
error $\leq10^{-8}$. The last successful analysis is selected before
examining outcome category. A successful earlier table cannot rescue
a missing or unfaithful final table.

The replay changes only imports/data paths, preserves frozen numeric
functions, and omits dispatch code. Original function text and
source/export digests are included. It reads bound archived analysis
values; participant programs remain unexecuted text. It reproduces the
eligibility test without inferring or regenerating tables.
\section{Narrow runtime proof audit}
\label{app:runtime-proof}

This appendix states a sufficient numerical contract for the continuous
projection argument and documents which parts were established locally.
The real-arithmetic argument, the archived recurrence computation, the
current machine implementation and the historical execution environment
are distinct evidence layers. None of the checks below adds a live agent
trial or reruns a physical episode.

\paragraph{Identity and trust assumptions.}
The inspected executable matches the experimental Python digest. It is
CPython 3.12.14 on ARM64. The current math image is
\texttt{libsystem\_m.dylib}; its UUID, text/constants hashes and relevant
function bytes were recorded. The current thread reports a zero floating-point
control register: round to nearest, with no flush-to-zero or alternative
handling enabled. We assume the ordinary IEEE binary64 semantics of
the decoded add, multiply, divide, fused multiply-add and square-root
instructions. The source-to-instruction interpretation is a manual
derivation with exact constant checks, not a proof-assistant-certified
binary analysis. The historical inventories pin Python but do not attest
the loaded math image or control register; equivalence to the historical
environment remains an explicit assumption.

\paragraph{Bound for the captured norm implementation.}
Let $u=2^{-53}$ and $d=2^{-1075}$, the unit roundoff and half the smallest
subnormal. Here $u$ denotes roundoff, not a physical action. For finite
binary64 coordinates, put $a=\max(|x_1|,|x_2|)$ and
$b=\min(|x_1|,|x_2|)$, and let $r=\sqrt{a^2+b^2}$ in exact arithmetic.
The captured function returns $a$ when $b=0$. Otherwise it drops $b$
only if the stored exponent difference exceeds 27. This implies
$b/a<2^{-27}$, including a subnormal $b$, so dropping it reduces the
exact norm by less than $r2^{-55}=ru/4$.

For nonzero retained inputs, the decoded branches use the following
scales, where $e_a$ is the biased exponent field of $a$:
\begin{center}
\begin{tabular}{@{}lll@{}}
\toprule
Exponent field & Input scale & Undo scale\\
\midrule
$e_a<539$ & $2^{768}$ & $2^{-768}$\\
$539\leq e_a<1534$ & $1$ & $1$\\
$e_a\geq1534$ & $2^{-768}$ & $2^{768}$\\
\bottomrule
\end{tabular}
\end{center}
The retained inputs scale exactly. In the small-input branch even the
smallest subnormal becomes $2^{-306}$. In the large-input branch a
nonzero retained $b$ has exponent difference at most 27 and remains
normal after scaling. In all branches the squared larger input and
the sum are normal and cannot overflow. The core computes
$t=\operatorname{RN}(a_s^2)$,
$v=\operatorname{RN}(b_s^2+t)$ by a fused operation, and
$w=\operatorname{RN}(\sqrt v)$, before undoing the scale.
Thus $w$ lies within factors $(1-u)^2$ and $(1+u)^2$ of the scaled norm.
Undoing a power-of-two scale is exact for a finite normal output;
subnormal rounding contributes at most $d$. Overflowed outputs are not
finite accepted norms. Combining the discarded-input allowance gives
the sufficient bound, for a finite returned norm $h$,
\begin{equation}
 r(1-3u)-d\ \leq\ h\ \leq\ r(1+3u)+d.
 \label{eq:hypot-bound}
\end{equation}
The rational checker verifies
$(1-u/4)(1-u)^2\geq1-3u$ and $(1+u)^2\leq1+3u$.
In particular, $h\leq4$ implies $r<5$ and
$|h-r|<H=16u$. This is an analytic bound for the captured sequence
under the stated instruction semantics, not an accuracy inference from
test samples.

\paragraph{Action normalization.}
Let $L$ be the exact binary64 value of $0.6$. An accepted but unscaled
action satisfies $h\leq L$, hence
$r\leq(L+d)/(1-3u)<L+8u$. When $h>L$, the validator rescales by
$aL/h$. Here $r>1/2$ and $h>1/2$. The compiled multiplication by a real
scalar and division by a positive real denominator each round each
component once; the zero imaginary terms do not introduce a separate
nonzero error. Using $\sqrt2<3/2$ gives
\begin{equation}
 \|a_{\rm applied}\|_2
 \leq \frac{L(1+u)^2}{1-3u-2d}
       +3d(1+u)+\tfrac32 d
 < L+8u.
 \label{eq:action-bound}
\end{equation}
This covers accepted rescaling, including subnormal allowances, rather
than treating the rejection tolerance as an applied-action allowance.

\paragraph{Two remaining library accuracy contracts.}
Write $\widehat s=\operatorname{pow}(h,2)$ and
$C(\phi)=\operatorname{cmath.exp}(i\phi)$ for the computed values.
The needed contracts are
\begin{equation}
 |\widehat s-h^2|\leq\epsilon_{\rm pow},\qquad
 \|C(\phi)\|_2\leq1+\delta_{\rm rot}.
 \label{eq:remaining-contracts}
\end{equation}
They must hold uniformly over relevant norm inputs and permitted finite
phases. The second includes the wrapper's \texttt{exp} and combined
\texttt{\_\_sincos\_stret} computation. No bound on angular error is
needed: every such control vector has norm at most
$(1+\delta_{\rm rot})(L+8u)$, irrespective of its direction.
Inspection of call targets establishes the operation path but does not
prove either accuracy contract.

\paragraph{Composed one-step bound.}
For a state before computed escape, $r<5$. Let $\alpha,\beta,c,L$
denote the exact binary64 task constants, so the reference drift is
$(\alpha+\beta r^2)x+c$. Set $B=32u+d$. Under the sufficient choices
$\epsilon_{\rm pow}=\delta_{\rm rot}=2^{-30}$, the relevant basic-operation
intermediates have magnitude below 32; the square is handled separately
by its library contract. Each basic rounding error is therefore at most
$B$. The squared-radius error is bounded by
$10H+H^2+\epsilon_{\rm pow}$. Two basic operations form the drift
multiplier. Scaling the state and adding the bias give vector error
at most $5\beta(10H+H^2+\epsilon_{\rm pow})+13B$.
The compiled complex product uses one ordinary product and one fused
operation per component; its error contributes less than $3B$.
Final vector addition contributes less than $3B/2$.

Consequently the sufficient loss allowance beyond the ideal $-L$ term is
\begin{equation}
 E=5\beta(10H+H^2+\epsilon_{\rm pow})
    +18B+8u+\delta_{\rm rot}(L+8u).
 \label{eq:composed-error}
\end{equation}
The exact rational evaluation is approximately
$6.2870891\times10^{-10}<10^{-9}$. For $\|q\|_2\leq1$ and a positive
projection lower bound $p$, this yields
\[
 q^\top x'\geq(\alpha+\beta p^2)p+q^\top c-L-E.
\]
The archived subtraction and downward floor therefore remain sufficient
\emph{if} the stated instruction, library and historical-equivalence
premises hold. The prior audit's minimum terminal gap exceeds $0.0371$;
it also clears the norm-test error. The same ten unique archived
projection certificates are covered conditionally, with no relabeling.

\paragraph{Checks and open obligations.}
The archived local audit checked exact inequality constants, captured function
identity and historical inventory coverage. The reader bundle rechecks only
the exact constants and exponent ranges with \texttt{artifact/numeric\_constants.py};
it does not recapture or attest the historical runtime. Fifty-two norm inputs,
including exponent/underflow edges, twenty square inputs and twenty-six
phases were probed in the archived audit interpreter. Every finite tested output
satisfies the sufficient bounds; an overflowed norm is recorded separately.
These are scalar diagnostics, not uniform accuracy proofs. Negative
controls reject altered function bytes and oversized error budgets, and
retain an open status when a library premise or historical binding is
missing. They test the checker, not the missing premises themselves.

The remaining obligations are explicit: prove both bounds in
Eq.~\eqref{eq:remaining-contracts} for the relevant library implementation,
and justify its equivalence and floating-point settings in the historical
runs. Current image hashes cannot supply absent historical attestation.
Accordingly, this addendum narrows the numerical assumptions and supports
their composition; it does not turn the archived results into unconditional
machine-verified claims. Exact route dynamic programming and the
observed trial counts are unchanged.


\section{The 32-trial common-menu protocol and panel usage}
\label{app:common-menu-protocol}

\paragraph{Selected histories and menus.}
Table~\ref{tab:common-menu-histories} identifies the four purposeful
histories using the earlier transfer panel's source IDs. The two
continuous histories are at physical step 10; the two route histories
are at step 2. Each appears with both profiles, both repeats and both
conditions. These are selected reference histories, not a random sample
of histories or eight newly sampled worlds.

\begin{table}[htbp]
\centering\small
\begin{tabular}{lllrrr}\toprule
Delivered ID & Source ID & Family & Instance & Cutoff & Menu size\\\midrule
probe-01 & snapshot-03 & Continuous & 1 & 10 & 13\\
probe-02 & snapshot-07 & Continuous & 2 & 10 & 13\\
probe-03 & snapshot-02 & Route & 1 & 2 & 3\\
probe-04 & snapshot-06 & Route & 2 & 2 & 3\\\bottomrule
\end{tabular}
\caption{Fixed histories for the 32-trial panel. Source histories and
menu identities are hash-bound in the archived allocation.}
\label{tab:common-menu-histories}
\end{table}

Both conditions receive identical history and task/menu bytes for a
given probe; the condition instruction differs. The continuous menu
is the twelve radius-$0.6$ directions at angles $2\pi j/12$,
$j=0,\ldots,11$, rounded to twelve decimals, followed by zero.
The route menu is $[(0,0),(1,0),(2,0)]$ in that order. Only the indexed
menu action is executed. The original physical clock, remaining
horizon and task objective are unchanged. The participant receives
visible history, public task/menu and its own analysis outputs, without
hidden dynamics or future observations. Trial order was fixed before
the run and is distributed in \texttt{artifact/data/crossing.json}.

\paragraph{Common output scaffold and condition-specific execution.}
The shared structured-output scaffold wraps a \texttt{decision} object
that can request analysis or return an action ID, a record, or both.
It does not itself enforce condition-specific fields, record identity,
dimensions, loss range or action agreement; the local contract does.
A record has exactly six fields: \texttt{schema}, \texttt{case\_id},
\texttt{menu\_sha256}, \texttt{loss\_definition\_id},
\texttt{candidate\_losses} and \texttt{rule}. It contains 1--128 rows,
one numeric loss in $[0,1]$ per menu action, and no weights. Each row
uses the same declared meaning: an uncalibrated forecast of failing the
original full task after that next action and the best continuation the
participant can plan. The two rules minimize the uniform row mean or
row maximum. Decimal numeric tokens are evaluated as exact rationals;
ties select the lowest menu index. Single and duplicate rows are legal,
and duplicates retain multiplicity. Rows are not verified world models.

In the optional condition, a legal \texttt{action\_id} is required and
an accompanying record is optional. Record failure or disagreement is
logged without replacing that action. In the mandatory condition, only
the record is allowed; its declared rule deterministically selects the
executed action. An invalid record produces no action and no retry.
Thus mandatory agreement is by construction, while optional agreement
is an observational check. This six-field contract differs from the
48-trial panel's weighted, four-field last-analysis adapter in Appendix E;
their eligibility counts are not measurements of the same interface.

\paragraph{Budgets and recorded use.}
For the 32-trial panel, each fresh trial permits up to three optional
local analyses and one final decision: at most four model requests,
3,780 seconds per trial and 900 seconds per request. Analysis has a
30-second CPU/wall limit, sampled 1 GiB RSS stop and 100,000-character
JSON output cap. It yields no new physical observation. No retry,
repair prompt, neutral request or outcome-dependent extension is
allowed; the panel ceiling is 128 model requests. The completed panel
used 48 analysis requests and 32 final decisions. Its single selected
action per trial was graded; crossed labels reuse the bound menu census,
not additional live observations.

\begin{table}[htbp]
\centering\small
\begin{tabular}{lrrrrr}\toprule
Panel & Requests & Input & Output & Cached input & Unknown calls\\\midrule
48-trial transfer & 120 & 2,325,280 & 484,134 & 372,352 & 0\\
32-trial common menu & 80 & 1,733,978 & 422,952 & 263,168 & 0\\\bottomrule
\end{tabular}
\caption{Separate participant-model usage ledgers. Cached input is already
included in input, and reasoning tokens, when reported, are part of output.
Neither category is added a second time. Pilot usage is in Appendix D.}
\label{tab:panel-usage}
\end{table}

Table~\ref{tab:panel-usage} reports archived receipt-derived counts, not
analyst usage or provider charges; these latter quantities were not
measured. Synthetic transport and privileged local computation are outside
the table. The 32-trial plan reserves 1,050,000 input plus 128,000 output
tokens for each request lacking usage, separately from measured totals.
Its final unknown-call count is zero, so no such reserve remains in that
panel's final accounting; this policy is not a provider billing cap.
The source-bound scalar projection is
\texttt{artifact/data/protocol-summary.json}. It provides provenance and
protocol identities without redistributing the full provider logs.


\paragraph{Exact requests and executable reader path.}
\texttt{artifact/protocol/} contains both condition instructions, the shared
wire schema, the four initial packets and all 80 request payloads, including
each trial's accumulated analysis records. Its checker renders every full
prompt and verifies the archived prompt and payload hashes. For example,
\begin{verbatim}
python3.12 -I -B -S artifact/protocol/check.py --show-call 1
\end{verbatim}
prints the first request. The historical collection entry point is supplied
as inspection-only text; provider transport credentials and full provider
responses are not redistributed. Configured profile names are not independent
attestation of server weights.

\texttt{artifact/menu\_replay.py} binds each public packet to its transfer
history and menu, then recomputes the 26 continuous and six route menu labels.
The original record comparison and this numerical replay are separate checks.
Continuous loss labels retain the open uniform-library and historical-runtime
premises in Appendix F; present-host agreement does not discharge them.

\section{Study inventory and historical boundaries}
\label{app:inventory}
The eight-trial pilot, 48-trial prospective transfer panel, 32-trial
common-menu panel, separate 48-assignment follow-up and new 16-target
prospective choice-prediction pilot are five evidence sets in the main paper. The two 48-slot studies are distinct. Their counts are
never pooled. The follow-up has 46 completed trials and two technical losses,
not 48 completed scientific trials. The two 1,600-draw development samples and the
5,000-draw main acquisition study appear only in the companion; the earlier
prediction criteria failed and are retained there. The acquisition result
is not an additional live LLM or crossing experiment.

Older archive labels describe different work: \emph{H4} denotes an unfinished
collection whose 84 remaining assignments stayed paused and whose crossing
status is unassessed due to missing evidence. The \emph{historical technical
null} denotes an earlier infrastructure-limited record, not a successful
zero-use scientific trial; its unresolved usage reservation stays outside
the pilot accounting. These records supply no extra evaluable trials here.
The archive's v1.8/v1.9 labels identify preserved predecessor research
versions, not additional studies or independent replications. Current paper
revision numbers track presentation and reader-access changes. This revision
additionally incorporates the prospective choice-prediction pilot, with its
original interrupted segment and separately completed last target block.
The earlier follow-up, predecessor scientific inputs and denominators, and
the primary acquisition procedure remain unchanged. The world-prediction
research design is unrun and contributes no result to this manuscript.

\section{Analysis-interface errors and consumed slots}
\label{app:interface-audit}
The historical shared instruction describes \texttt{history.states} and
\texttt{history.actions}; the evaluator actually supplies a Python dictionary.
Five analysis calls failed on \texttt{history.states}. Two other calls failed
because \texttt{next} is absent from the restricted built-ins; the historical
instruction refers to the existing interface without listing those built-ins.
These are observed interface-use errors, not missing final decisions. The
seven affected trials are listed below. Each made three analysis attempts:
one failed and two succeeded. Errors consumed an attempt and were visible
to the next request. They did not trigger replacement trials or extra slots.

\begin{table}[htbp]\centering\small
\begin{tabular}{lllll}\toprule
History & Profile & Condition & Repeat & Error\\\midrule
probe-01 & gpt-6-astra & Mandatory & 1 & Dictionary attribute\\
probe-01 & gpt-6-astra & Optional & 2 & Dictionary attribute\\
probe-01 & gpt-6.1-sol & Mandatory & 2 & Dictionary attribute\\
probe-01 & gpt-6.1-sol & Optional & 1 & Unavailable \texttt{next}\\
probe-01 & gpt-6.1-sol & Optional & 2 & Unavailable \texttt{next}\\
probe-02 & gpt-6-astra & Optional & 1 & Dictionary attribute\\
probe-02 & gpt-6.1-sol & Optional & 2 & Dictionary attribute\\\bottomrule
\end{tabular}
\caption{All seven affected common-menu trials. All are continuous-task
trials; both profiles requested max reasoning effort. The first row is the
sole rule-sensitive mandatory record. These are descriptive associations,
not estimates of the errors' effects on actions or records.}
\end{table}

The 32 common-menu trials produced 48 analyses: 41 successful and seven
erroneous. In the separate transfer panel, all 24 route trials made zero
analysis attempts, while the 24 continuous trials made 72 attempts: 67
successful and five erroneous. These denominators are not pooled. The
frozen adapter still supports only one transfer record and no complete pair.

\texttt{artifact/interface\_audit.py} reconstructs all 32 trial rows from the
last cumulative analysis list in each trial's archived requests, checks every
earlier list as a prefix, and reconciles attempt counts with the final trial
result and remaining-slot field. It also accounts for the 48 transfer trials.
No archived submitted program is executed. This procedure avoids counting
the same error again each time it appears in a later prompt.

The reconstructed tables and their crossing identities remain checkable.
However, the observed scarcity and row structure occurred under this
imperfect interface. A corrected description could change analysis use,
record structure or actions; no such behavioral effect has been measured.
Local checks of dictionary access and built-in availability cannot answer
that causal question. A prospective comparison must preserve old records
and compare the two interface descriptions under a separately specified
allocation and budget. Mandatory reporting and external execution would
still be a combined intervention in either interface arm.

\paragraph{Pilot sensitivity reader boundary.}
\texttt{artifact/pilot\_sensitivity.py} regenerates all twelve single-candidate
plans and their five selection metrics from the disclosed component inputs,
then reconstructs the five selected indices and native-action compatibility.
It now re-evaluates all twelve first actions and compares the full archived
witness/certificate records as described in Appendix~\ref{app:b-completions}.
This extends the reader through the conditional task labels. The numerical
and historical-runtime premises of the grade evidence remain open.


\section{Follow-up contract, accounting and verification boundary}
\label{app:followup}
The follow-up profiles are gpt-6-astra/max/Fast and gpt-6.1-sol/max/Fast.
Each has 24 original assignments, eight per arm and eight matched triplets.
The same four historical contexts are used, with two repeats per profile;
neither repeats nor individual requests are new independent task samples.
A and B execute a legal proposal. A may omit its record; an invalid optional
record does not invalidate an otherwise legal proposal. B and C require a
valid record and legal proposal; a protocol-invalid final produces no action.
C executes the declared record selector and retains the proposal. The common
schema checks case/menu identity, dimensions, finite values and legal actions.
Single rows, duplicates and same-action rules are allowed. The shared
analysis interface uses dictionary access and states its available operations,
addressing the earlier interface mismatch; this is not a randomized old/new
interface comparison, so its causal effect is unmeasured.

Each trial permits at most three optional analyses and one final decision,
four model requests, 900 seconds per request and 3,780 seconds per trial.
Errors consume a slot; model responses are not repaired and failed
assignments are not replaced. The closed results record 36 analyses for
Astra and 36 for Sol, with zero recorded analysis errors in each. These counts
do not turn the historical interface comparison into a controlled experiment.
Decisions were bound before privileged label lookup. Follow-up labels reuse
the same identity-bound finite menu census, retaining the continuous numerical
and historical-runtime premises of the earlier paper. No new physics or
live crossed model decision was generated for this analysis.

\paragraph{Three execution segments and missingness.}
The original 48-position allocation was retained across three separately
authorized execution segments. Positions 1--12 completed; 13 and 14 ended
with HTTP 429 technical losses in the first and second segments. The third
segment completed positions 15--48 without repeating either failed position.
Both technical losses are Astra, in A and C of the same matched triplet.
Thus Astra has 22 completed, two technical and zero protocol losses; Sol has
24 completed and zero technical or protocol losses. All 48 assignments started
and have result records; none remains active or unstarted. Completion of the
34-position last segment does not mean all 48 assignments scientifically
completed. Missing pairs are excluded only from observed-pair descriptions,
with their denominators visible; no missing-at-random claim is made.
The final segment's 88 request receipts contain no terminal 429 diagnostic.
That observation does not identify the earlier limiter or establish a
permanent resolution; credit exhaustion is not inferred.

\begin{table}[htbp]
\centering\small
\begin{tabular}{lrrrrr}\toprule
Profile & Requests & Known/unknown & Input & Cached input & Output\\\midrule
Astra & 60 & 58/2 & 1,428,745 & 311,168 & 305,022\\
Sol & 60 & 60/0 & 1,418,971 & 263,424 & 294,997\\\bottomrule
\end{tabular}
\caption{Follow-up accounting only. Token counts cover known usage.
Cached input is included in input; reasoning output is included in output.
Neither subset is added again. Requests are not independent trials.}
\label{tab:followup-usage}
\end{table}

Using the frozen tariffs, measured credit equivalents are 1,336.90190 for
Astra and 264.37032 for Sol. Astra's two unknown requests retain a separate
1,690 credit-equivalent reserve, based on 1,050,000 input plus 128,000 output
tokens per request. These are local accounting quantities, not a provider
invoice or provider token ceiling. The original 7,000 local admission bound
was not increased. Measured usage and unknown reserves remain separate;
cache splits are measured for known receipts, while any cold-price treatment
of an unknown split would be an assumption, not measured zero cache.

\paragraph{Schema denominators.}
The retained field \texttt{partial\_triplets} counts planned blocks missing
an A, B or C assignment; its zero value does not assert scientific completion.
The reader separately reports fully completed triplets (Astra 7, Sol 8),
scientifically incomplete triplets (Astra 1, Sol 0), and triplets containing
technical losses (Astra 1, Sol 0). The retained field
\texttt{cache\_split\_unknown\_requests} counts unknown cache splits only
among requests with known token usage: its assessed denominators are
58 for Astra and 60 for Sol. Astra's two usage-unknown requests are excluded
from that count and retain their full separate reserves. Thus zero unknown
cache splits is neither zero cached input nor known usage for every request.
The machine-readable \texttt{artifact/followup/SCHEMA.json} defines these
populations and exclusions explicitly.

\paragraph{Reader inputs and remaining scope.}
The package contains 46 exact minimal finals, including 31 original numeric
tables and declared rules, plus the common/A/B/C instructions and wire schema.
The reader checks trial, profile, block, case and menu identities; reparses
decimal tokens as exact rationals; recomputes rule selections and row diagnostics;
and reconciles decisions, commit hashes and aggregates. It checks A/B against
the proposals and C against the record selector. All 46 completed labels are
re-read from the bound \texttt{artifact/data/crossing.json}; none changed.
Section~\ref{sec:reader-access} gives the source paths and executable entry point. No participant
analysis or task physics is run by this source-table check.

The visible follow-up chain is in \texttt{artifact/followup/interactions/}.
Its index binds each call to its trial and segment. Exact prompts and response
files retain their bytes; analysis request files contain the submitted code
and public history, and tool-output files contain the archived JSON result.
All visible agent-message text is supplied in order; the last message is the
selected terminal response. Earlier visible messages are not forwarded as
chat history: only the selected analysis program and its result enter the
next request's \texttt{analysis\_records}. The reader verifies this distinction
from the cumulative prompts. This is 120 calls, not 177 calls or 177 trials.

Intent and receipt projections omit paths and account/thread identifiers.
Raw provider streams, including hidden reasoning, remain local; their source
hashes establish a binding but do not let readers reconstruct omitted fields.
Fifty-three collector and dependency files across the three segments are
inspection-only \texttt{.py.txt} files, with source and distributed hashes and
explicit path-redaction status. They document collection without providing
a runnable redispatch command. The reader's corruption controls reject
missing calls, changed identities, wrong conditions, broken code/result/history
links and altered finals. These are tests of the supplied evidence chain,
not new scientific trials or execution of participant programs.

Full operational archives remain local. Hashes do not establish independent
laboratory replication or immutable provider-model/tier identity. The proposed population,
live query-selection and representation-transfer studies were not run;
the finite-acquisition companion retains its separate, unchanged sample and claims.


\subsection{History-level outcome accounting}
\label{app:followup-history}
Table~\ref{tab:followup-histories} retains the four purposeful histories and
their two repeats within each profile. It describes these assignments only;
neither the cells nor their repeated contexts are independent new worlds.
The reader derives this table and the main-text family, transition and
rule-sensitivity tables from the bound minimal finals, allocation and menu
census. Technical assignments retain no scientific label. All completed
follow-up labels are P or L; there are no U or protocol outcomes.
% Generated from identity-bound existing follow-up sources.
\begin{table}[htbp]
\centering\small
\begin{tabular}{llllll}\toprule
Profile & History & Family & A: P/L/T & B: P/L/T & C: P/L/T\\\midrule
Astra & probe-01 & Continuous & 1/1/0 & 2/0/0 & 1/1/0\\
Astra & probe-02 & Continuous & 2/0/0 & 2/0/0 & 1/1/0\\
Astra & probe-03 & Route & 2/0/0 & 2/0/0 & 2/0/0\\
Astra & probe-04 & Route & 1/0/1 & 2/0/0 & 1/0/1\\
Sol & probe-01 & Continuous & 1/1/0 & 0/2/0 & 2/0/0\\
Sol & probe-02 & Continuous & 2/0/0 & 2/0/0 & 1/1/0\\
Sol & probe-03 & Route & 2/0/0 & 2/0/0 & 2/0/0\\
Sol & probe-04 & Route & 2/0/0 & 2/0/0 & 2/0/0\\
\bottomrule
\end{tabular}
\caption{All follow-up history/arm outcomes; T denotes technical loss with no scientific label. Every cell has two planned assignments. All completed labels are P or L. Repeated contexts are not independent draws of new worlds; P is continuation feasibility, not measured full-task completion.}
\label{tab:followup-histories}
\end{table}


\section{Supporting fixed-contract control}
\label{app:structural-control}

The comparison holds the archived tables, rule definitions, menu and task-label
map fixed. It removes operations rather than withholding information from a
baseline. We evaluate (i) native actions and rule equality, (ii) native actions
and the single-distinct-row condition, and (iii) independent per-record
evaluation of each supported rule. The last baseline never needs a partner's
table. Only after producing its results do we compare them with paired cells.

\paragraph{Proposition (per-record sufficiency under a fixed contract).}
For fixed tables $A_i$, a common menu, and a fixed family $\mathcal R$ of
deterministic rules with specified tie breaking, the collection
$\{(i,r,r(A_i)):r\in\mathcal R\}$ determines every supported-rule-family
cell and every observed-rule cell whose rule belongs to $\mathcal R$.
\emph{Proof.} Each cell is one of those evaluations; no partner table is an
argument of $r(A_i)$. A fixed action-label map also determines its label.
This statement does not establish the validity of that map or identify
an actor's internal computation.

Two weaker checks need fewer rule evaluations. If $r_1=r_2$, then
$c_{i1}=c_{i2}=a_i$. If a table has a single distinct row, mean and worst case
select the same column. Duplicate-containing tables with other distinct
rows do not satisfy this condition. The full-scan equality is thus a
compositional identity of the contract. The
empirical question is how often the weaker checks suffice and what contrasts
remain; agreement with the full scan is not an independent discovery of the
identity. We measure recoverable cells, not a wall-clock speedup or causal
attribution accuracy.


\begin{table}[h]
\centering\small
\begin{tabular}{p{0.52\linewidth}rr}
\toprule
Baseline information & Observed cells & Family cells\\
\midrule
Native actions + identical observed rules & 32/32 & ---\\
Native actions + one-distinct-row check & --- & 29/32\\
Independent two-rule scan per record & 32/32 & 32/32\\
Additional cells from full table after scan & 0 & 0\\
\bottomrule
\end{tabular}
\caption{Mandatory condition only. Counts describe recoverable displayed
cells, not independent observations or accuracy on external diagnoses.}
\end{table}


\section{Prospective choice-prediction protocol and source access}
\label{app:prospective}
\paragraph{History construction and target branches.}
The preparation fixed two instances per disclosed task family with seed
2026103001, before model outputs. The purpose was a small feasibility panel
with fresh histories and legal continuation opportunities, rather than a
population sample or a balanced difficulty comparison. In continuous
control, a privileged preparation controller generated six transitions by
choosing the menu action with the smallest next-state norm. In route tasks,
the fixed action prefix was $(0,0),(1,0),(2,0),(0,0)$; initial fuel was set to
retain a known feasible continuation at the source and both prescribed
branches, with a three-unit margin. Those preparation controllers are not
participant-model behavior. All twelve source/target packets passed the
precollection continuation check; the continuous preparation would stop on a
failed witness rather than resample. The four histories differ from the
earlier published histories. The original construction code and generation
record are supplied for inspection in \texttt{artifact/prospective-design/};
the reader does not execute them.

\begin{table}[htbp]
\centering\small
\begin{tabular}{llcl}\toprule
Histories & Family & Source $\to$ target step & Two separate branch actions\\\midrule
h01, h02 & Continuous & $6\to7$ & $(0.6,0)$; $(-0.6,0)$\\
h03, h04 & Route & $4\to5$ & $(0,0)$; $(2,0)$: advance 1; 3 nodes\\
\bottomrule\end{tabular}
\caption{Fixed source--target construction. Step counts completed transitions.
Each target preserves the entire source history and adds one transition and
its observation under the indicated exogenous action. The two branches are
alternatives from the same source, not consecutive transitions. Branch actions
were fixed before source-model outputs and need not equal the source proposal.}
\label{tab:prospective-branches}
\end{table}


The model actors receive the generated public histories, not the privileged
preparation state or continuation witnesses. Each target contains the full
source prefix and one additional action--observation transition. The target
branches and source-proposed action are separate quantities.

\paragraph{Allocation and information boundaries.}
The pilot fixes four new purposeful histories in the two existing disclosed
families, two profiles, one source per history/profile and two precommitted
target branches per source: eight sources, 32 observers and 16 targets.
Eight initial references make 64 assigned roles. Each target branch is a
precommitted branch from the source history; it need not follow the model's
source-proposed action. Both observers receive that fact. Every role uses a
fresh context; the source never receives future target snapshots. The source's
private analysis transcript is not replayed to observers or target, although
its information can be summarized in the record. The four prose fields and
numeric table are passed only to the with-record observer. The target receives
neither record nor observer responses. Exact role instructions, schemas,
packets, allocation, prompts and selected terminal responses are supplied.

The source final schema requires the action, numeric table, declared rule
and four prose fields; the target final schema requires the action alone.
Both observer arms receive both schemas and complete actor instructions.
This evaluates prediction across these output protocols; a same-protocol
target arm and an isolated protocol-change effect are outside the allocation.
Source and target actors may request up to three local analyses, followed by
a final choice; observers make one request without an analysis tool. Both
observer arms have the same requested model/max/Fast profile and 900-second
request allowance. Actor allowance is 3,780 seconds, each local analysis
30 seconds. These allowances do not require equal realized token use, and
the added record increases input length. Requested profiles and fresh local
contexts are recorded; immutable remote model identity is not attested.

\paragraph{Reference gate and completion.}
Two hand-authored policy cases are evaluated with and without a record for
each profile. The initial gate requires eight valid transport/protocol results
and all four with-record predictions correct. Astra's with/without accuracy
is 2/2 and 2/2; Sol's is 2/2 and 1/2. Without-record accuracy is descriptive,
not a gate. The later Sol completion repeats these four references: all valid,
with-record 2/2 and without-record 1/2. Only then are the two remaining
predictions locked and the target requested. Reference success does not
establish scientific prediction benefit or durable provider availability.

The original segment has 61 completed, one interrupted and two unstarted
roles. The new segment has seven completed roles, including four new controls
and the three remaining scientific roles. Their union covers all 64 original
assignments, but the 68 completed segment roles are not 68 independent samples.
The interrupted original request is neither deleted nor marked successful.
Completion was decided after partial-result inspection and adds no new
history. It restores the planned target denominator without transforming this
pilot into a confirmatory sample. The original partial panel remains visible.

\paragraph{Scoring and structural diagnostics.}
For each valid target, an observer scores one iff its action ID equals the
target's menu-bound ID. Invalid observer protocol would score zero; unavailable
source or target leaves the paired difference in $[-1,1]$. The primary
descriptive quantity averages the two target differences within each history,
then the four histories, separately by profile. There are no protocol losses
in the completed original or completion results. We report assigned, started,
completed, interrupted and unstarted roles separately. Empty denominators
would be unavailable, never success. There is no pooled primary estimate,
population significance test or model ordering.

The reader derives source validity, declared selection, both supported-rule
choices, row count, distinct-row count and proposal match directly from the
eight exact finals. It checks all menu/action identities. All eight source
records match; seven are single-row and one Sol record has two distinct rows.
All eight select the same action under both rules, forced in the seven
single-row cases and observed in the remaining one. No row duplication occurs.
The source-action copy baseline is legal for every target. The numeric table
plus prose is one treatment. The semantic informativeness of the prose was
not independently evaluated, and the seven single-row tables do not establish
that the prose is uninformative. Table~\ref{tab:prospective-paired} reports
joint correctness for the combined panel; equal marginal counts can include
different correctly predicted targets. Identical-input target variability
and task-specific informative-record sensitivity were not separately measured.
The family and history breakdowns in
Tables~\ref{tab:prospective-families} and~\ref{tab:prospective-histories} were
added after inspecting the combined results. They preserve every assigned
history and target and leave the prespecified aggregate unchanged.

\begin{table}[htbp]
\centering\small
\begin{tabular}{lllrrr}\toprule
Profile & History & Family & With record & Without record & Net correct\\\midrule
Astra & h01 & Continuous & 1/2 & 2/2 & $-1$\\
Astra & h02 & Continuous & 1/2 & 1/2 & $0$\\
Astra & h03 & Route & 2/2 & 2/2 & $0$\\
Astra & h04 & Route & 2/2 & 2/2 & $0$\\
Sol & h01 & Continuous & 2/2 & 1/2 & $+1$\\
Sol & h02 & Continuous & 1/2 & 1/2 & $0$\\
Sol & h03 & Route & 0/2 & 0/2 & $0$\\
Sol & h04 & Route & 0/2 & 1/2 & $-1$\\
\bottomrule\end{tabular}
\caption{Post hoc breakdown by the eight source-history/profile units.
Each row contains its two linked target branches. Sol h04 includes the
separately completed final target; the original fifteen-target panel remains
in Table~\ref{tab:prospective}. Histories are not repeated independent draws.}
\label{tab:prospective-histories}
\end{table}


\paragraph{Usage and reserves.}
The initial segment has 96 requests: 95 measured and one unknown Sol request.
Astra's 47 measured requests used 747,683 input and 282,909 output tokens,
including 150,016 cached input tokens; Sol's 48 measured requests used 854,673
input and 345,366 output tokens, including 84,480 cached input tokens.
The new segment uses seven Sol requests: four controls, two observers and one
target, with 88,601 input and 34,066 output tokens, including 32,000 cached
input tokens. Cache is a subset of input, and reasoning is a subset of output;
neither is counted again. The new segment contains no analysis request.

Frozen local credit equivalents are 1,013.60680 for initial Astra,
250.12470 for initial Sol, and 22.85310 for the Sol completion. The original
unknown request still reserves 1,050,000 input plus 128,000 output tokens,
equivalent to 169 credits. Across segments: 103 issued requests, 102 measured,
one unknown; measured equivalents 1,286.58460, reserve-inclusive 1,455.58460.
The initial 10,000 local bound remains; the separately approved completion
bound is 1,690 with at most ten new requests. Neither is a provider invoice
cap. Analyst use and older follow-up reserves are outside this accounting.
All 102 measured receipts have known cache splits; this does not make the
unresolved request's usage known. No 429 appears in these receipts; the
original interruption's cause remains undetermined.

\paragraph{Prospective evidence and check limits.}
The prospective evidence retains its own exact manifest within the single
supplementary archive. It includes
103 exact prompts and their schemas, 102 selected terminal responses,
34 analysis program/result pairs, eight source records, and projected request
receipts and journal chronology. Analysis code is readable data; the reader
never executes it or a physical simulator. The main package binds the
supplement's manifest. Its source reader reconstructs prompts, re-parses
source tables, checks source--prediction--target order, and recomputes initial
and combined panels and usage. The implementation shares the collector's
parser and scalar readout; negative controls challenge its evidence links.
Local source hashes and projected timestamps are not independent provider
attestation. Raw account identifiers, hidden reasoning and provider-added
context are not supplied. The reader audits the supplied evidence, not
unobserved internal computation or general-world prediction.


\section{Evidence access and operational accounting}
\label{app:reader-access}
Table~\ref{tab:evidence-access} specifies access separately for each evidence
set. The bundle contains LaTeX sources, disclosed current-study mechanisms,
older numerical replay and the follow-up source reader. The older panel's
80 payloads reconstruct prompt hashes; archived participant programs are
never executed. The unchanged companion is optional supporting material. Its catalog,
ordered draw, 5,000 histories, queries, certificates and replay are in a
separate companion bundle; they are not needed for the main-paper checks. Existing-input replay is not new data collection or
historical process attestation.

For the follow-up, the standard-library reader parses exact numeric lexemes
from 46 minimal finals, including 31 tables. It derives validity, both rule
choices and row diagnostics, compares archived decisions, verifies menu and
profile identities, and checks 46 labels by lookup in the bound census.
A separate interaction reader now reconstructs all 120 archived application
prompts from the condition instructions, four fixed public packets and the
exact cumulative analysis prefix. It links 118 terminal response files,
all 177 visible agent messages and 72 analysis input/output pairs to the
same allocations and finals. Two failed requests retain no response and
unknown usage. Collector sources are supplied as inspection-only text;
participant programs and collector code are never executed by these checks.
Intent and receipt files are explicit field projections with original source
hashes. This permits an audit of the supplied application-visible chain,
not independent attestation of historical execution, provider-added context
or hidden reasoning. Original streams with machine/account identifiers remain
local. The earlier revision added numerical re-evaluation of the existing
pilot sensitivity actions. The paper also reports the
prospective choice-prediction pilot and its separately recorded completion.
The prospective supplement contains all eight exact source records, 32
scientific predictions, 16 target finals and 12 reference finals; 103 exact
prompts, 102 terminal response files and 34 analysis code/result pairs support
reconstruction of the visible inputs. The unresolved request has a prompt and
an unknown-usage reserve. Its reader checks action/menu identities and
chronology and recomputes both result panels without executing submitted
programs or physics. It shares the collector's record parser and scalar
calculation, so this is reproducibility, not an independent implementation.
The prospective component is bound by hash to the main component. Both are restored from the single supplementary archive.

\begin{table}[htbp]
\centering\small
\begin{tabular}{>{\raggedright\arraybackslash}p{0.19\linewidth}>{\raggedright\arraybackslash}p{0.36\linewidth}>{\raggedright\arraybackslash}p{0.36\linewidth}}\toprule
Evidence & Supplied material & Boundary of the check\\\midrule
Prospective supplement & Eight exact source records; 103 prompts, 102 terminal responses, 34 analysis pairs; allocation, menus, instructions and projected receipts & Recompute original/combined predictions and record structure; retain interrupted attempt; projected chronology is not provider attestation\\
Follow-up sources & 46 minimal finals, including 31 numeric tables and rules;
allocation, common/A/B/C instructions, schema, menu and bound labels & Recompute
final decisions and labels by lookup; no provider or historical-runtime attestation\\
Follow-up interaction & 120 prompts, 118 terminal responses, 177 visible messages, 72 analysis input/output pairs; inspection-only collector sources &
Reconstruct cumulative application prompts and bind finals; projected receipts do not independently authenticate historical execution or provider context\\
Older common-menu panel & All 80 request payloads, prompt renderer, instructions,
histories and label-replay inputs & Reconstruct archived prompts and hashes;
participant programs are inspected as data, never executed\\
Pilot selector sensitivity & Twelve original action grades with full witnesses
and certificates, action lock, fixed plan inputs and reader & Recompute plans,
selectors and conditional labels; compare bound evidence for every action\\
Mechanisms and labels & Disclosed current-study mechanisms and older numerical
replay sources & Conditional continuous premises remain open;
replay is not independent historical or laboratory attestation\\\bottomrule
\end{tabular}
\caption{Reader access by evidence set. Final-table recomputation, visible
interaction consistency and historical attestation are separate capabilities.
No hidden model reasoning is included.}
\label{tab:evidence-access}
\end{table}


The single supplementary archive stores repeated UTF-8 lines once in readable
JSON shards, with an indexed path-to-line mapping and the original file
hashes. PDF files remain directly readable. The extraction utility rejects
unexpected paths, invalid line references and altered bytes before restoring
the two complete components. Component manifests and the main-to-prospective
binding are then verified by the evidence reader. This representation changes
storage only; it preserves all source files and both execution segments.
\noindent\begin{minipage}{\linewidth}
To obtain the original directory layout for direct inspection, run:
\begin{verbatim}
python3.12 -I -B -S verify.py --extract-to ../ruu-evidence
\end{verbatim}
\end{minipage}

\noindent\begin{minipage}{\linewidth}
From
\texttt{../ruu-evidence/paper-package}, the component reader is:
\begin{verbatim}
python3.12 -I -B -S verify_all.py --prospective-dir ../prospective
\end{verbatim}
\end{minipage}

The additional scalar check derives the A2 strict-dominance comparisons,
the four-cell paired table, and the family and history breakdowns directly
from the supplied numeric source and terminal responses. It also checks
the fixed source--target prefixes and branch actions. It does not execute
participant programs or the history generator.

\end{document}
